\chapter{Pearson Advanced Electrochemistry: Complete Exercise Bank}
\label{chap:pearson_all_questions}

\begin{problembox}
\textbf{Problem 1} \hfill \textbf{[Pearson Advanced Electrochemistry Q12]}
\label{q:ext-pearson-1}

A metal having negative reduction
of three metallic cations,X,Y and Z are
potential, when dipped in the solution of
+0.52, -3.03 and -1.18 V, respectively.
its own ions,has a tendency to
The order of reducing power of the
(a)remain as metal atoms
corresponding metals is
(b)become electrically positive
(a)Y>Z>x
(c) become electrically negative
(b)x>Y>Z
(d)be deposited from the solution
(c) Z>Y>x
(d) Z>X>Y
13.The calomel electrode is reversible with
9.A gas'X'at 1 atm is bubbled through
respect to
a solution containing a mixture of 1 M
(a) Hg+
-Yand1M-Zat25$/circ$C.If the reduction
(b)H+
potential of Z>Y>X, then,
(c)Hg$2$+
(a)Y will oxidize X and not Z
(d) Cr
(b)Y will oxidize Z and not X
14.Which one of the following does not get
(c)Y will oxidize both X and Z
oxidized by bromine water?
(d)Y will reduce both X and Z
(a)Fe$2$+ to Fe3+
10.S
Standardelectrodepotentialdataare
(b) Cutto Cu$2$2+
useful for understanding the suitability
(c)Mn$2$+ to MnO4
of an oxidant in a redox titration.Some
(d) Sn2+to Sn4+
half-cellreactions andtheir standard
potentials are givenbelow:
\end{problembox}

\begin{problembox}
\textbf{Problem 2} \hfill \textbf{[Pearson Advanced Electrochemistry Q15]}
\label{q:ext-pearson-2}

In salt bridge, normally KCl is used
because
MnO(aq)+ 8H"(aq) + 5e$2$ $/to$ Mn2+
(aq)+4HO()E$/circ$=1.51 V
(a)it is a strong electrolyte.
(b) it is good conductor of electricity.
CrO$2$(aq) + 14H*(aq)+ 6e $/to$ 2Cr+
(c)K+and Cl ions have nearly same
(aq)+7HO(l)E$/circ$=1.38 V
ionic mobility.
Fe$2$(aq)+e$/to$Fe$2$+(aq)E$/circ$=0.77 V
(d) it is an ionic compound.
Cl(g)+2e$/to$2CI(aq)E$/circ$=1.40 V
\end{problembox}

\begin{problembox}
\textbf{Problem 3} \hfill \textbf{[Pearson Advanced Electrochemistry Q16]}
\label{q:ext-pearson-3}

By how much would the oxidizing power
of MnO/Mn$2$+ couple change if the H+
Identifytheonlyincorrectstatement
ions concentration is decreased 100 times
regarding the quantitative estimation of
at 25$/circ$C?
aqueous Fe(NO)
(a) increases by 189 mV
(a)MnO can be used in aqueous HCI
(b) decreases by 189 mV
(b) Cr,O can be used in aqueous HC1
(c)will increase by 19 mV
(c)MnO can be used in aqueous HSO
(d)will decrease by 19 mV
(d) Cr,O$2$- can be used in aqueous HSO4
17.The solution of CuSO, in which copper
\end{problembox}

\begin{problembox}
\textbf{Problem 4} \hfill \textbf{[Pearson Advanced Electrochemistry Q11]}
\label{q:ext-pearson-4}

Thedecreasingorderofstandardelectrode
rod is immersed, is diluted to 10 times.
potential of Mg, K, Ba and Ca, is
Thereductionelectrodepotential
(a)K, Ca,Ba, Mg
(a) increases by 0.0295 V
(b)Ba,Ca,K, Mg
(b) decreases by 0.0295 V
(c)Ca,Mg,K,Ba
(c) increases by 0.059 V
(d)Mg,Ca,Ba,K
(d) decreases by 0.059 V
--- PAGE 3 ---
Electrochemistry8.3
\end{problembox}

\begin{problembox}
\textbf{Problem 5} \hfill \textbf{[Pearson Advanced Electrochemistry Q18]}
\label{q:ext-pearson-5}

The standard reduction potential of
24.The standardreduction potentials at25$/circ$C
oxygen in acidic solution is +1.23 V. What
of Li'|Li, Ba$2$|Ba, Na |Na and Mg$2$Mg
is the standard reduction potential of
are -3.05,-2.73,-2.71 and -2.37 V,
oxygen in basic solution?
respectively. Which is the strongest
(a)+0.404 V
(b)-0.404 V
reducing agent?
(c) +2.056 V
(d)-2.056 V
(a) Li
(b)Ba
19.The standard reduction potentials of
(c) Na
(d) Mg
Cu$2$+|Cu and Cu$2$+|Cu* are 0.337V
Some standard electrode potentials are
and 0.153 V,respectively. The standard
given:
electrode potential of Cu*|Cu half-cell is
Fe$2$++ 2e$/to$Fe; E$/circ$=-0.440 V
(a) 0.184 V
Fe3++3e$/to$Fe;E$/circ$=-0.036V
(b) 0.827V
(c) 0.521 V
The standard electrode potential for:
(d) 0.490 V
Fe3++e$/to$Fe$2$, is
\end{problembox}

\begin{problembox}
\textbf{Problem 6} \hfill \textbf{[Pearson Advanced Electrochemistry Q20]}
\label{q:ext-pearson-6}

The electrode potential of hydrogen
(a)-0.476V
(b)-0.404 V
(c) +0.988 V
electrode in neutral solution and 298K is
(d)+0.772 V
(a) -0.413 V
(b)zero
\end{problembox}

\begin{problembox}
\textbf{Problem 7} \hfill \textbf{[Pearson Advanced Electrochemistry Q26]}
\label{q:ext-pearson-7}

The dissociation constant for CH,COOH
(c)-0.826V
(d) +0.413 V
is 1.8 $/times$ 10- at 298 K.The electrode
\end{problembox}

\begin{problembox}
\textbf{Problem 8} \hfill \textbf{[Pearson Advanced Electrochemistry Q21]}
\label{q:ext-pearson-8}

Electrode potential will be more for
potential for the half-cell: Pt  H2
hydrogen electrode at pH (at the same
(1 bar)|0.5 M-CH,COOH,at298 K
temperature)
is (log 2 = 0.3; 1og 3 = 0.48; 2.303RT/F
(b)3
= 0.06)
(a) 4
(c) 2
(d) 5
(a)-0.3024 V
(b) -0.1512 V
22.S
Saturated solution of KNO,is used to
(c)+0.3024 V
make'salt bridge'because
(d) +0.1512 V
(a)velocity of K+is greater than that of
NO:
At 25$/circ$C,the solubility product of CuCl
(b) velocity of NO;is greater than that
of K+
The value of Ecucu
is (log 2 = 0.3;
(c) velocities of both K* and NO are
2.303 RT/F= 0.06)
nearly the same
(d)KNO is highly soluble in water
(a) -0.274 V
(b) -0.402 V
(c) +0.53 V
(d)+0.402V
23.The standard reduction
potentials
of PtCr,O$2$, Cr+3;PtMnO,Mn+2;
\end{problembox}

\begin{problembox}
\textbf{Problem 9} \hfill \textbf{[Pearson Advanced Electrochemistry Q28]}
\label{q:ext-pearson-9}

The standard reduction potential of
Pt|Ce+4, Ce+ in the presence of acid are
oxygen in acidic solution is 1.23 V (O2
1.33 V,1.51V and 1.61V, respectively,at
+4H,O*+ 4e$/to$6HO).The standard
25$/circ$C.The decreasing order of oxidizing
reduction potential of oxygen in basic
poweris
solution is (2.303RT/F=0.06)
(a) CrO$2$->MnO>Ce+4
(a) -1.23 V
(b) MnO>CrO$2$->Ce+4
(b)-0.39 V
(c)Ce+>MnO>CrO$2$
(c) +0.39 V
(d)MnO>Ce4>CrO$2$
(d) +2.07 V
--- PAGE 4 ---
8.4Chapter8
29.The standard potentials of MnO|Mn$2$+
30.The following reactions represent the
and MnO2| Mn$2$ electrodes in acid
reduction of IO,ion into I ion in acidic
solution are 1.51 and 1.23 V, respectively.
andbasic medium.Predictinwhich medium
Standardelectrodepotentialfor
IOion will act as a better oxidizing agent?
the electrode,MnO4|MnO2in acid
IO+6H++6e$/to$I+3H2O; 
solutionis
E$/circ$=+0.907 V
(a) +1.697 V
IO+ 3HO+ 6e$/to$I+ 6OH;
(b) +5.09 V
E$/circ$=+0.260V
(c) +0.28 V
(a) Acid medium
(b)Basic medium
(d) +1.37 V
(c)Equally in both
(d)Notpredictable
GalvanicCell
31.The correct cell diagram for the following
\end{problembox}

\begin{problembox}
\textbf{Problem 10} \hfill \textbf{[Pearson Advanced Electrochemistry Q34]}
\label{q:ext-pearson-10}

The cell reaction for the given cell is
reaction and E$/circ$for the cell is
spontaneous if Pt,Cl (P, atm)|CI |Cl
2AgBr(s)+H(g)$/to$ 2Ag(s)+ 2H*+ 2Br
(P,atm), Pt
E$/circ$AgBrAsB-= +0.10 V.
(a) P>P
(a)(Pt) H| H+II Br"|AgBr |Br (Pt);
(b) P<P
E$/circ$=0.10V
(c) P,=P2
(b) (Pt) H|H+II Br"|AgBr |Br (Pt);
(d) P,= 1 atm
E$/circ$=-0.10V
(c)(Pt) Br,|AgBr |Br I H+|H (Pt);
Which one of the following statements
E$/circ$=0.10V
is incorrect regarding an electrochemical
(d) (Pt) Br,| AgBr | Br  H+|H (Pt);
E$/circ$=-0.10V
cell?
(a) The electrode on which oxidation
32.In an experimental
1set-upforthe
takes place is called anode.
measurement of EMF of a half-cellusinga
reference electrode and a salt bridge,when
(b)Anode is the negative pole.
the salt bridge is removed, the voltage
(c)The direction of the current is same
(a)remains the same
as that of the direction of flow of
(b)increases to maximum
electrons.
(c)decreases half the value
(d) The flow of current is partly due to
(d) drops to zero
flow of electrons and partly due to
33.After some time,the voltage of an
flow of ions.
electrochemical cell becomes zero.This is
because
36.When an electric current is drawn from a
(a) their electrode potential becomes
galvanic cell
zero.
(a)EMF suddenly increases.
(b) their reduction potential become
(b)EMF gradually increases and attains
equal but have opposite sign.
a maximum value.
(c) their reduction potential become
equal and have the same sign.
(c) EMF decreases and finally falls to
(d)theions of the electrolyte in the salt
zero.
bridge stop moving.
(d)EMF must remain constant.
--- PAGE 5 ---
Electrochemistry8.5
37.Identification of anode and cathode in an
44.Use of lithium metal as an electrode in
electrochemical cell is made by the use of
high energy density batteries is due to
(a) Galvanometer
(a)lithium is the lightest element.
(b) Salt bridge
(b)lithium has the highest oxidation
(c)Voltmeter
potential.
(d)Potentiometer
(c)lithium is quite reactive.
38.For the cell Zn|Zn$2$I Cu$2$|Cu, if the
(d) lithium does not corrode readily.
concentration of both,Zn$2$+and Cu$2$ions
45.A depolarizer used in dry cell batteries is
are doubled,the EMF of the cell
(a) NHCI
(b)MnO2
(a) doubles
(c)KOH
(d) Na,PO4
(b) reduces to half
46.Which is correct about fuel cells?
(c)remains same
(d) becomes zero
(a) Cellcontinuously run as long as fuels
are supplied.
39.The standard EMF of a galvanic cell can
(b)These aremore efficient and freefrom
be calculated from
pollution.
(a)the size of the electrode
(c)These are used to provide power and
(b)the pH of the solution
drinking water to astronauts in space
(c)the amount of metalin the anode
programme.
(d)the E$/circ$values of the half-cells
(d) Allof these
47.When a lead storage battery is discharged
\end{problembox}

\begin{problembox}
\textbf{Problem 11} \hfill \textbf{[Pearson Advanced Electrochemistry Q40]}
\label{q:ext-pearson-11}

The value of equilibrium constant for a
feasible cell reactionmust be
(a) SOis evolved
(b)lead sulphate is consumed
(a) <1
(b)Zero
(c) lead is formed
(c) =1
(d) >1
(d) sulphuric acid is consumed
41.If the cell reaction is spontaneous then
48.For a cellreactioninvolving atwo-
(a) $/Delta$G$/circ$=-ve
electron change,the standard EMF of
(b) E$/circ$red = -ve
the cell is found to be 0.295 V at 25$/circ$C.
(c) E$/circ$e=+ve
The equilibrium constant of thereaction
(d)$/Delta$G=-ve
at 25$/circ$C will be
(a) 1$/times$1010
(b) 1$/times$10-10
42.Which of the following pair of metals,
(c) 29.5$/times$ 10-2
(d) 2$/times$1010
when coupled,will give maximum EMF
49.The Ecen for Ag(s)|AgI (satd)IlAg*
for a voltaic cell?
(0.10 M)|Ag (s) is +0.413 V. What is the
(a)Fe and Cu
value of K of AgI?
(b) Pb and Au
(c) Cu and Au
(a) 1.0$/times$ 10-8
(b) 1.0$/times$10-7
(c) 1.0 $/times$10-14
(d) 1.0$/times$10-16
(d) Ca and Cu
50.Assuming that hydrogen behaves as an
43.When a lead storage battery is charged
ideal gas, what is the EMF of the cell at
(a)PbO dissolves.
25$/circ$C if P = 600 mm and P = 420 mm:
(b) the lead electrode becomes coated
Pt|H (P)|HCI|H(P)|Pt? [Given:
with lead sulphate.
2.303RT/F= 0.06, log 7 =0.85]
(c)sulphuric acid is regenerated.
(a)-0.0045 V
(b)-0.0V
(d) the amount of acid decreases.
(c)+0.0045 V
(d) +0.0015 V
--- PAGE 6 ---
8.6Chapter8
51.For the electrochemical cell,M|M* I
57.F
For a reaction A(s)+ 2B* $/to$A$2$++ B(s); Kc
X |X, EsM= 0.44 V and Exx-=0.33 V.
has been found tobe 10.The EMF of
From this data,one can deduce that
the cell is
(a) M + X $/to$ M*+ Xis the spontaneous
(a) 0.354V
(b) 0.708 V
reaction
(c) 0.534V
(d) 0.453 V
(b) M++ X- $/to$ M + X is the spontaneous
\end{problembox}

\begin{problembox}
\textbf{Problem 12} \hfill \textbf{[Pearson Advanced Electrochemistry Q58]}
\label{q:ext-pearson-12}

An electrochemical cell is set up as
reaction
follows:
(c) Ec=0.77 V
(d) E=-0.77 V
Pt|H(s) (1 atm)|0.001 M HCl II
0.1 M HA |H(s) (1 atm) |Pt
\end{problembox}

\begin{problembox}
\textbf{Problem 13} \hfill \textbf{[Pearson Advanced Electrochemistry Q52]}
\label{q:ext-pearson-13}

The EMF of the cell:Zn|Zn$2$+(0.01 M)Il
EMF of this cell is zero because (pKa
Fe$2$+ (0.001 M)|Fe at 298 K is 0.2905 V,
of HA = 5)
thenthevalue of equilibrium constant
(a) molar concentrations of acids are
for the cell reaction is
different
0.32
0.32
(a)e0.0295
(b) 100.0295
(b)Temperature is constant
0.26
(c)pH of two solutions are same
0.32
(c)100.0295
(d)100.0591
(d)both are standard hydrogen electrodes
\end{problembox}

\begin{problembox}
\textbf{Problem 14} \hfill \textbf{[Pearson Advanced Electrochemistry Q59]}
\label{q:ext-pearson-14}

.The
53.For the cell reaction:4Br +O2+ 4H*
reaction:Hz(g)+2AgCl(s)
$/to$ 2H*(aq) + 2CI(aq) + 2Ag(s) occurs in
2Br,+ 2HO; E$/circ$= 0.18 V.The value of
the galvanic cell
(log K) at 298 K is [2.303RT/F = 0.06]
(a)Ag|AgCl (s)|KCl(aq)|AgNO(aq)|
(a) 12
(b) 6
(c) 18
(d) 3
Ag
(b)Pt|H(g)|HCI(aq)|AgNO(aq)|Ag
54.The standard EMF for the cell reaction:
(c)Pt|H(g)|HCI(aq)|AgCl(s)| Ag
Zn(s) + Cu$2$+(aq)$/to$Zn$2$(aq)+ Cu(s) is
(d) Pt |H2(g)|KCl(aq)|AgCl(s)|Ag
1.10 volts at 25$/circ$C.The EMF of the cell
60.The standard reduction potentials in
reaction when 0.1 M Cu2+ and 0.1 M Zn2+
acidic conditions are 0.77V and 0.53 V,
solutions are used at 25$/circ$C is
respectively, for Fe$2$| Fe2t and I | I
(a) 1.10 v
(b) 1.041 V
couples. The equilibrium constant for
(c) -1.10 V
(d) -1.041 V
the reaction: 2Fe$2$++312Fe$2$+I,is
55.For the cell: Ni | Ni$2$+Il Cu$2$*| Cu;
(2.303RT/F = 0.06)
E$/circ$=0.77 V. By which of the following
(a) 2$/times$108
(b) 108
(c) 10*
(d) 10-8
activity,E will increase?
(a) On decreasing [Nit2]
61.The following electrochemical cell has
(b) On decreasing [Cu+]
been set up: Pt(s)Fe$2$,Fe$2$+(a =1)Il Ce,
(c) On increasing mass of Nielectrode
Ce$2$+(a=1)|Pt(s);E$/circ$(Fe$2$Fe$2$)=0.77 V;
(d) On increasing mass of Cu electrode
E$/circ$(CeICe$2$)=1.61V.If an ammeter
is connectedbetween the two platinum
56.The standardEMFof aDanielcellat298
electrodes, predict the direction of flow
Kis E,.When theconcentration of ZnSO
of current.Will the current increase or
is 1.0 M and that of CuSOis 0.01 M, the
decrease with time?
EMF becomes E, at 298 K.The correct
(a)Ce electrode to Fe electrode,decrease
relationship between E, and Eis
(b)Ce electrode toFeelectrode,increase
(a)E,=E
(b) E=0
(c)Fe electrode to Ce electrode,decrease
(c) E,>E
(d)E<E
(d)Fe electrode to Ce electrode,increase
--- PAGE 7 ---
Electrochemistry8.7
62.For the reaction:H(1 bar)+ 2AgCl(s)
\end{problembox}

\begin{problembox}
\textbf{Problem 15} \hfill \textbf{[Pearson Advanced Electrochemistry Q64]}
\label{q:ext-pearson-15}

ThestandardpotentialsofOCl/Cl and
2Ag(S)+ 2H+(0.1 M)+2CI (0.1 M);
CI/Clare 0.94V and-1.36V,respectively.
G$/circ$=-48,250Jat 25$/circ$C.The EMFof cell
The E$/circ$value of OCI/Cl,will be
in which the given reaction takes place is
(a)3.24 V
(b)-0.42 V
(a) 0.25 V
(c)-2.30 V
(d) 0.52v
(b) 0.37V
65.From the following E$/circ$ values for the
(c) 0.13 V
half-cells:
(d) 0.0.49 V
(i) D $/to$D$2$+ 2e'; E$/circ$=-1.5 V
63.A cell contains two hydrogen electrodes.
The negative electrode is in contact
(ii)B++e$/to$B;E$/circ$=-0.5V
with a solution of 10-$/circ$M hydrogen
(i)A$2$-$/to$A$2$+e;E$/circ$=1.5V
ions.The EMF of the cell is 0.118 V at
(iv) C$2$++e$/to$C*;E$/circ$=+0.5 V
25$/circ$C.The concentration of hydrogen
ions at the positive electrodeis
Which combination of two half-cells
(a) 10- M
would result in a cell with largest
(b) 10-$2$ M
potential?
(c) 10-4 M
(a) iandii
(b) iand iv
(d) 10-s M
(c) ii and iv
(d) i and iv
Electrolysis
66.Which of the following statements does
69.If mercury is used as a cathode during the
notdifferentiatebetweenelectrochemical
electrolysis of an aqueous NaClsolution,
cell and electrolytic cell?
the ions discharged at cathode are
(a)Spontaneousor non-spontaneous
(a) H*
nature of the chemical process.
(b) Na*
(b) Chemical reactions occurring at the
(c)OH
electrodes
(d) CI
(c)Positive and negative nature of anode.
(d)EMF measurement.
67.The electrode through which electrons
70.Electrochemical equivalent is more for
enter theelectrolyticsolutionis
(a)Hydrogen
(a) cathode
(b) Silver
(b) anode
(c) Copper
(c)may be anode or cathode
(d)Zinc
(d) both,anode and cathode
Which process occurs in the electrolysis
71.The electrolytic bath used in gold plating
of an aqueous solution ofnickel chloride
of copper articles contains
at nickel anode?
(a) Molten gold
(a) Ni $/to$ Ni2+ + 2e"
(b) Copper sulphate(aq)
(b) Ni$2$+ + 2e$2$ $/to$ Ni
(c)2CI $/to$ Cl+ 2e"
(c) AuCl,(aq)
(d) 2H*+2e$/to$H
(d) AuCl, + NaCN(aq)
--- PAGE 8 ---
8.8Chapter8
\end{problembox}

\begin{problembox}
\textbf{Problem 16} \hfill \textbf{[Pearson Advanced Electrochemistry Q72]}
\label{q:ext-pearson-16}

Copper can be deposited from acidified
77.A certain current liberated 0.50 g of
copper sulphate and alkaline cuprous
hydrogen in 2 h. How many grams of
cyanide. If the same current is passed for
copper can be liberated by the same
a definite time
current flowing for the same time in a
(a) the amount of copper deposited from
copper sulphate solution?(Cu=63.5)
acidic copper sulphate will be higher
(a) 12.7 g
(b) 15.88 g
(b) the amount of copper deposited
(c) 31.75 g
(d) 63.5 g
from alkaline cuprous cyanide will be
78.A current of 3.7 A is passed for 6 h
higher
(c) the same amount of copper will be
between nickel electrodes in 0.501 of 2 M
solution of Ni(NO).The molarity of
deposited
Ni$2$+ at the end of electrolysis is
(d) copper will not deposit in either case
(a) 1.172 M
(b) 0.172 M
73.In the electrolytic cell, flow of electrons is
(c) 0.586M
(d) 2 M
from
79.The
current
efficiency
of
an
(a) Cathode to anode in solution
electrodeposition of copper metal in
(b) Cathode to anode through external
which 9.8 g of copper is deposited by a
supply
current of 3A for10000 s,from aqueous
(c) Cathode to anode through internal
copper sulphate solution,is about
supply
(a) 60/%
(b) 99/%
(d)Anode to cathode through internal
(c) 92/%
(d)75/%
supply
80.On passing electricity through dilute
74.Electrolytic cellis used to convert
H,SO4 solution, the mass of substances
(a) Chemical energy to electrical energy
liberated at the cathode and anode are in
(b)Electricalenergy to chemical energy
the ratio of
(c)Chemical
1energytomechanical
(a) 1:8
(b) 8:1
energy
(c) 1:32
(d) 1:16
(d)Electrical
energy
tomechanical
81.The electrochemical equivalents of two
energy
substances are E,and E.The current that
must pass to deposit the same amount at
75.Faraday's law of electrolysis fails when
the cathodes in the same time must be in
(a) temperature isincreased
theratio of
(b)inert electrodes are used
(a)E:E
(b)E:E
(c)a mixture of electrolytes is used
(c)(E-E):E
(d)E:(E-E)
(d)in none of these cases
82.T
The same quantity of electricityispassed
76.Using same quantity of current,which
throughonemolar solutionofHSO4
among Na, Mg and Alis deposited more
and one molar solution of HCl.The
(by mass)during electrolysis of their
amount of hydrogen evolved from H,SO
molten salts?
as compared to that from HClis
(a) Na
(a) the same
(b) Mg
(b) twice as such
(c)A1
(c) one half as such
(d) All in same
(d)dependent onsize ofelectrode
--- PAGE 9 ---
Electrochemistry8.9
83.In the electrolysis of acidified AgNO;
series.The mass ratio of iron deposited at
solution using Pt-electrodes, the anode
cathodes in the two cells will be
reaction is
(a) 3:1
(b) 2:3
(a) 2NO$/to$2NO2+O2+2e
(c) 1:1
(d) 3:2
(b)NO$/to$NO+O2+e"
88.A
A galvanic cell is set up from a zinc bar
2
(c) 2HO$/to$4H*+O2+4e"
weighing 100 gand 1.0 L of 1.0 M copper
(d) Pt$/to$Pt$2$++3e"
sulphate solution. How long would the
cell run if it is assumed to deliver a steady
\end{problembox}

\begin{problembox}
\textbf{Problem 17} \hfill \textbf{[Pearson Advanced Electrochemistry Q84]}
\label{q:ext-pearson-17}

Twoplatinum electrodeswereimmersed
current of 1.0 A? (Zn = 65.4)
in a solution of CuSO and electric
(a)53.6h
(b) 26.8 h
current was passed through the solution.
(c) 81.97 h
(d) 40.99 h
After some time,it was found that colour
of CuSOdisappearedwith theevolution
89.An ion is a reduced to the element when
of gas at the electrode. The colourless
it absorbs 6x 10"electrons.The number
solution contains
of equivalents of the ion is
(a)platinum sulphate
(a) 0.10
(b) 0.01
(b) copper sulphate
(c) 0.001
(d)0.0001
(c)copper hydroxide
90.In the lead storage battery, the anode
(d) sulphuric acid
reaction is Pb(s)+ HSO4+ HO
A solution containing 1.0 M each of
$/to$PbSO4(s) + HO*+ 2e".How many
Cu(NO),Mg(NO),AgNO, ,Hg(NO)
grams of Pb will be used up to deliver 1 A
isbeingelectrolysed usinginertelectrodes.
for 100 h? (Pb =208)
Thevaluesofstandardelectrodepotential
(a) 776 g
(b) 388g
are: Ag|Ag= 0.80 V, Hg$2$+|Hg =0.79 V,
(c) 194 g
(d) 0.1 g
Cu$2$+|Cu=0.34V,Mg$2$+| Mg=-2.37 V.
With increasing voltage,the sequence of
91.The copper anode of a cell containing
silver nitrate solution weighs 60.0 g.
deposit of metals on the cathode will be
Afterpassing currentfor some time,itis
(a)Ag, Hg,Cu, Mg
found that 3.24 g of silver if deposited on
(b)Mg,Cu, Hg,Ag
the platinum cathode.What is the final
(c)Ag, Mg, Cu
weight of the anode?(Ag = 108, Cu =64)
(d) Cu, Hg,Ag
(a) 0.96 g
(b) 60 g
86.1
During theelectrolysis ofanaqueous
(c) 59.04 g
(d) 60.96 g
salt solution, the pH in the space near
92.The quantity of electricity required for
one of the electrode was increased and
the reduction of 1 mole of FeOto Fe is
the other one was decreased.The salt
(a) 1F
(b) 0.33F
solutionwas
(c)3F
(d) 6F
(a)NaCl(very dilute)
(b)ZnCl,
93.Three faradays of electricity is passed
(c) NaCi(Conc.)
through molten AlO3, aqueous solutions
(d) Cu(NO)
ofCuSOandmoltenNaCl.The
amounts of Al, Cu and Na deposited at
87.Two electrolytic cells, one containing
the cathodes will be in the molar ratio of
acidified ferrous chloride and another
(a) 1:2:3
(b) 3:2:1
acidifiedferricchloride,areconnectedin
(c) 1:1.5:3
(d) 6:3:2
--- PAGE 10 ---
8.10Chapter8
94.Electrolysis of a solution of HSO4ions
101.If a current of 1.0 A is drawn from the
produces SOg.Assuming 75/% current
Daniel cell for 96.5 min, the cathode will
efficiency,what current should be
gain in weight by (Cu = 63.5, Zn = 65.4)
employed to achieve a production rate of
(a) 1.905 g
(b) 1.962 g
1 mole of SO$2$per hour?
(c) 3.81g
(d) 3.924 g
(a) 71.5 A
(b) 35.7 A
102.The current required to produce oxygen
(c) 53.0 A
(d) 143 A
at the rate of 2.8 ml per second during
95.The number of Faradays required to
electrolysis of acidulated wateris
produce1g-atom of Mgfrom MgClis
(a) 48.25 A/s
(a) 1
(b) 2
(b) 24.12 A/s
(c) 0.5
(d) 4
(c) 96.5 A/s
\end{problembox}

\begin{problembox}
\textbf{Problem 18} \hfill \textbf{[Pearson Advanced Electrochemistry Q96]}
\label{q:ext-pearson-18}

Passage of 96,500 coulomb of electricity
(d) 0.0048 A/s
liberates
L of Oat 273$/circ$C and
103.Sodium
amalgam
1isprepared
1by
2 atm during electrolysis.
electrolysis of aqueous NaCl using 10 g
(a) 5.6
(b) 16.8
mercury as cathode. How many Faraday
(c) 22.4
(d) 11.2
of electricity is required to prepare 18.7/%
Na-amalgam,by weight,with a current
97.An electrolytic cell contains a solution
efficiency of 50/%?
of AgSO4 and platinum electrodes.A
(a) 0.1 F
(b) 0.2F
current is passed until 1.6 g of O has
(c)0.05F
(d) 0.16 F
been liberated at anode. The amount of
silver deposited at cathodewould be
104.Element A (atomic mass = 112) and
(a) 108.0 g
(b) 1.6 g
element B(atomic mass = 27) form
(c) 10.8 g
(d) 21.6 g
chlorides.Solutions of these chlorides
98.When 12,000 coulombs of electricity
are electrolysed separately and it is found
that when the same quantity of electricity
is passed through the electrolyte, 3.0 g
is passed,5.6 g of A was deposited
of a metal of atomic mass 96.5 g/mol
while only 0.9 g of B was deposited. The
is deposited.The electro-valency of the
valency of B is 3. The valency of A is
metal cation in the electrolyte is
(a) 1
(b) 2
(a) +4
(b) +3
(c)3
(d) 4
(c) +2
(d) -4
99.How many electrons flow when a current of
\end{problembox}

\begin{problembox}
\textbf{Problem 19} \hfill \textbf{[Pearson Advanced Electrochemistry Q105]}
\label{q:ext-pearson-19}

The same current was passed successively
5 A is passed through a solution for 200 s?
throughsolutionofzinc-ammonium
sulphate andnickel-ammonium sulphate
(a)6.022$/times$1023
rendered alkaline with ammonia.The
(b) 6.24$/times$ 1021
weights of zinc and nickel deposited in
(c) 6.024 $/times$1021
a certain time were found to be 22.89 g
(d) 6.022$/times$1020
and 20.55 g,respectively. Given that the
\end{problembox}

\begin{problembox}
\textbf{Problem 20} \hfill \textbf{[Pearson Advanced Electrochemistry Q100]}
\label{q:ext-pearson-20}

The current of 9.65 A flowing for 10 min
chemicalequivalent weight of zinc is 32.7,
deposits 3.0 g of a metal. The equivalent
what is the chemical equivalent weight of
weight of the metalis
nickel?
(a) 10
(b) 30
(a) 58.71
(b) 29.36
(c) 50
(d) 96.5
(c) 14.39
(d) 36.42
--- PAGE 11 ---
Electrochemistry8.11
Conductance
106.Which of the following solutions have
\end{problembox}

\begin{problembox}
\textbf{Problem 21} \hfill \textbf{[Pearson Advanced Electrochemistry Q112]}
\label{q:ext-pearson-21}

The specific conductance of AgCl
highest resistance?
solutioninwaterwasdeterminedtobe
(a) 1N-NaC1
1.8 $/times$10-$/circ$Q- cm- at 298 K.The molar
(b) 0.05N-NaC1
conductances at infinite dilution, of Ag*
(c) 2N -NaC1
and Cl are 67.9 and 82.1 Q-' cm$2$mol-,
(d) 0.1N - NaC1
respectively. What is the solubility of
AgCl in water?
107.Variation of molar conductance of an
(a) 1.2 $/times$ 10-8 M
electrolyticsolutionwith temperatureis
(b) 1.44 $/times$ 10-1 M
that it
(a)increases withincrease of temperature
(c) 1.2 $/times$ 10- M
(b) decreases
increase
of
(d) 1.44 $/times$ 10-16 M
with
temperature
113.Equivalence conductance at infinite
(c)first increases then decreases
dilution of NHCl,NaOH and NaCl
(d)is not affected by temperature
are 129.8, 217.4 and 108.9 Q-1 cm$2$
108.Which pure substance will not conduct
mol',respectively. If the equivalent
electricity?
conductance of 0.01 N solution of
(a) Molten NaCl
NHOH is 9.532 Q-'cm$2$mol', then the
(b)Molten KOH
degree of dissociation of NHOH at this
(c) Liquefied HCI
temperatureis
(d) Liquid Hg
(a) 0.04/%
(b) 2.1/%
109.The correct order of molar conductance
(c) 4.0/%
(d) 44.7/%
at infinite dilution of LiCl, NaCl and
KCIis
\end{problembox}

\begin{problembox}
\textbf{Problem 22} \hfill \textbf{[Pearson Advanced Electrochemistry Q114]}
\label{q:ext-pearson-22}

The resistance of 1 M - CH,COOH
(a) LiCl>NaCl>KCI
solution is 250 Q,when measured in a
(b)KCl>NaCl>LiCl
cell of cell constant 125 m-'.The molar
(c)NaCI>KCI>LiCl
conductivity,in Q-m$2$mol-' is
(d) LiCI>KCI>NaCl
(a) 5.0 $/times$ 10-4
(b) 500
\end{problembox}

\begin{problembox}
\textbf{Problem 23} \hfill \textbf{[Pearson Advanced Electrochemistry Q110]}
\label{q:ext-pearson-23}

The molar conductance of a strong
(c) 2$/times$10-3
(d) 200
electrolyte atinfinite dilution
115.How does the electrical conductivity of
(a)tends to a finite value, which is above
20 ml of 0.2 M-MgSO4change when
that at higher concentration
0.5 M -Ba(OH), solution is gradually
(b) tends to a finite value,which is below
added in it, to excess?
that at higher concentration
(c) tends to zero
(a)decreases continuously
(b)increases continuously
(d) tends to a finite value,which is equal
(c)increases and then decreases
to that athigh concentration
(d) decreases and then increases
\end{problembox}

\begin{problembox}
\textbf{Problem 24} \hfill \textbf{[Pearson Advanced Electrochemistry Q111]}
\label{q:ext-pearson-24}

The best conductor of electricity is a
0.1 M solution of
\end{problembox}

\begin{problembox}
\textbf{Problem 25} \hfill \textbf{[Pearson Advanced Electrochemistry Q116]}
\label{q:ext-pearson-25}

The equivalent conductivity (in Q-1
(a) Boric acid
cm$2$ eq) of 1.0 M -HSO4 solution of
(b) Sulphuric acid
specific conductance 2.6$/times$ 10-  cm-", is
(c) Acetic acid
(a) 1.3$/times$ 102
(b) 6.5$/times$10l
(d)Propanoic acid
(c) 1.3 $/times$ 10-1
(d) 2.6$/times$102
--- PAGE 12 ---
8.12Chapter8
117.Themolarconductanceofa0.01M
122.Equal volumes of 0.015M-CHCOOH
solution of acetic acid was found to be
and 0.015 M -NaOH solutions are
16.30 Q-  cm-  mol' at 25$/circ$C.The ionic
mixed together.What would be the molar
conductances ofhydrogenand acetateions
conductivity of mixture if conductivity
at infinite dilution are 349.8 and 40.9 Q-
of CHCOONa is 6.3x10-4S cm-?
cm mol-', respectively, at the same
(a) 0.84 S cm$2$mol-1
temperature.What percentage of acetic
(b) 8.4 S cm$2$ mol-l
acidis dissociated at this concentration?
(c) 84 S cm$2$mol-1
(a) 0.04172/%
(b) 4.172/%
(d) 42 S cm$2$mol-1
(c) 41.72/%
(d) 0.4172/%
\end{problembox}

\begin{problembox}
\textbf{Problem 26} \hfill \textbf{[Pearson Advanced Electrochemistry Q123]}
\label{q:ext-pearson-26}

Calculate Am(in Q-  cm$2$mol')for SrCl
118.The distancebetween two electrodes of a
at 25$/circ$C,from the following data:
cell is 2.5 cm and area of each electrode is
5 cm$2$.The cell constant is
Conc.
0.25M
1.0 M
(a) 0.5 m-1
(b) 12.5 cm$2$
$/Delta$ (in Q- cm$2$mol-1)
260
250
(c) 2.0 cm
(d) 50 m-l
(a) 270
(b) 265
\end{problembox}

\begin{problembox}
\textbf{Problem 27} \hfill \textbf{[Pearson Advanced Electrochemistry Q119]}
\label{q:ext-pearson-27}

The molar conductivity of NHCl,OH"
(c) 240
(d) 275
and Cl at infinite dilution is 150,200
and 75 Q-' cm$2$mol-', respectively. If the
\end{problembox}

\begin{problembox}
\textbf{Problem 28} \hfill \textbf{[Pearson Advanced Electrochemistry Q124]}
\label{q:ext-pearson-28}

The resistance of a solution A is 50 Q and
molar conductivity of a 0.01 M-NHOH
that of solution Bis100 $/Omega$,both solutions
solution is 22Q-cm$2$mol-', then its
being taken in the same conductivity cell.
degree of dissociation is
If equal volumes of solution A and B
(a) 0.146
(b) 0.063
are mixed, what will be the resistance of
(c)0.080
(d) 0.293
the mixture using the same cell? Assume
that there is no increase in the degree of
\end{problembox}

\begin{problembox}
\textbf{Problem 29} \hfill \textbf{[Pearson Advanced Electrochemistry Q120]}
\label{q:ext-pearson-29}

Calculate the ionic product of water at
dissociation of A and B on mixing.
25$/circ$Cfrom the following data:
(a) 150 $/Omega$
Conductivity of water=5.5
(b)75$/Omega$
$/times$ 10-mho m-1
(c)33.33Q
*= 0.035 mho m$2$mol-1
(d) 66.67 $/Omega$
/textasciicircum()$/circ$oH = 0.020 mho m$2$mol-1
\end{problembox}

\begin{problembox}
\textbf{Problem 30} \hfill \textbf{[Pearson Advanced Electrochemistry Q125]}
\label{q:ext-pearson-30}

In a conductivity cell, the two platinum
(a)2$/times$10-14M$2$
electrodes, each of area 10 cm$2$ are fixed
(b)1$/times$10-$2$M$2$
1.5 cm apart. The cell contained 0.05 N
(c)1$/times$10-8M$2$
solution of a salt.If the two electrodes are
(d) 1 $/times$10-14 M2
just half dipped into the solution which
has a resistance of 50 Q, the equivalent
121.Calculate K. of acetic acid if its 0.05 M
conductance of the salt solution, in Q-
solution has molar conductivity of 7.814
cm$2$eq-', is
$/times$10-Q-m$2$molat25$/circ$C.Given:Afor
(a) 120
CH,C0OH=3.907$/times$10-$2$Q-m$2$mol"
(b) 60
(a) 2$/times$10-s
(b) 1.8 $/times$ 10-s
(c)240
(c) 4 $/times$10-4
(d) 0.02
(d) 3000
--- PAGE 13 ---
Electrochemistry8.13
AnswerKeys-ExerciseI
ElectrodePotential
\end{problembox}

\begin{problembox}
\textbf{Problem 31} \hfill \textbf{[Pearson Advanced Electrochemistry Q99]}
\label{q:ext-pearson-31}

(b) 100. (c) 101. (a) 102. (a) 103. (b) 104. (b)105. (b)
Conductance
106.(b) 107. (a) 108.(c) 109.(b) 110.(a) 111.(b) 112.(c) 113.(c) 114.(a)115.(d)
\end{problembox}

\begin{problembox}
\textbf{Problem 32} \hfill \textbf{[Pearson Advanced Electrochemistry Q116]}
\label{q:ext-pearson-32}

(a) 117. (b) 118. (d) 119. (c) 120. (b) 121. (a) 122. (c) 123. (a) 124. (d) 125. (a)
--- PAGE 14 ---
8.14Chapter8
EXERCISEII(JEEADVANCED)
SectionA(OnlyoneCorrect)
1.Astudentmadethefollowingobservations
4.The standard reduction potential for the
in the laboratory:
reactions:Ag*+e"$/to$Ag and Ag(NH)2
+e"$/to$Ag+ 2NH,are +0.79 V and +0.37 V,
(I) Clean copper metal did not react
respectively. From these values and the
with1 M-Pb(NO)solution.
Nernst equation,what should be K,for the
(II) Clean lead metal dissolves in 1 M
Ag(NH)*ion?[Given: 2.303RT/F =0.06]
AgNO,solution and crystals of
(a) 1.0 $/times$10-7
Ag metal appeared.
(b) 1.0$/times$10?
(III)Clean silver metal did not react with
(c) 2.15$/times$ 1019
1 M-Cu(NO)solution.
(d) 4.64 $/times$ 10-20
The order of decreasing
reducing
5.The overall formation constant for the
character of the threemetalsis
reaction of 6 mole of CN with cobalt (I1)
(a) Cu>Pb>Ag
is1x10'.What is theformationconstant
(b) Cu>Ag>Pb
forthereactionof6molesofCNwith
(c)Pb>Cu>Ag
cobalt (IIl)?Given that
(d) Pb>Ag>Cu
Co(CN)$2$+e$/to$Co(CN)-;
2.Wehave
anoxidation-reduction
E$/circ$=-0.83V
system:[Fe(CN)]$2$-+e"=[Fe(CN)];
Co$2$++e$2$$/to$Co$2$+;E$/circ$=+1.81V
E$/circ$=+0.36 V.The ratio of concentrations
of oxidized and reduced from at which
and 2.303RT/F=0.06.
thepotentialofthesystembecomes
(a) 1.0$/times$ 1063
(b) 1.0$/times$102s
0.24 V, is [Given: 2.303 RT/F = 0.06)
(c)1.0 $/times$ 10-25
(d) 1.0$/times$ 10-63
(a) 2:1
(b) 1:2
(c) 1:20
(d) 1:100
6.For the process:Cu$2$++ 2e"$/to$Cu;log[Cu$2$]
vs.Era graphis shown in the figure,where
3.The standard reduction potential for the
OA = 0.34 V. The electrode potential of
process:[Co(HO)]$2$++e$2$$/to$[Co(HO)]$2$+
the half-cell of CulCu$2$+ (0.1 M) will be
is1.8V.The standard reductionpotential
[2.303RT/F= 0.06]
fortheprocess:[Co(NH)]$2$+e
$/to$[Co(NH)]$2$+is 0.1V.Which of the
complexion,[Co(HO)$2$+or[Co(NH)]$2$+
can be oxidized to the corresponding
cobalt (II1) complex, by oxygen, in basic
A
medium,under standard condition?
(a) [Co(H,O)]$2$+
log[Cu2+]
(b) [Co(NH)]$2$+
(c)both
(a) -0.31 V
(b) +0.31 V
(d) none of these
(c)-0.37V
(d)+0.37 V
--- PAGE 15 ---
Electrochemistry8.15
7.Tell which of the following statements
12.The EMF of the cell: Hg(l)|HgCl,(s),
accuratelydescribestheeffectofaddingCN
KCl sol.(1.0N)|Quinohydrone|Pt,is
to the cathode of a cell with a cell reaction:
0.210 V at 298 K.What is the pH of the
Cd+2Ag*$/to$2Ag+Cd$2$,E$/circ$=1.2V
quinohydrone solution, the potential of
(a)E$/circ$ increases because Cd(CN)$2$-forms
the normal calomel electrode is 0.279 V
(b)E$/circ$decreases because Cd(CN)$2$forms
andE$/circ$for the quinohydrone electrode
(c)E$/circ$increases because Ag(CN),forms
is 0.699 V,both at the same temperature.
(d)E$/circ$ decreases because Ag(CN), forms
[2.303RT/F = 0.06]
(a) 3.5
(b) 7.0
8.For a electrochemical cell Zn|Zn$2$+(C
(c) 1.85
M) Cu$2$+(C M)| Cu, the decrease in
(d)-3.5
free energy at a given temperature is a
13.The potential (EMF)of a cell consisting
function of
of an anode of silver in 0.10 M-AgNO
(a) In C
(b) In C
solution and a cathode ofPtimmersed
(c) In CC,
(d) In C/C
in a solution of 1.6 M - Cr,O$2$,
9.Select the correct option if it is known
= 0.80 V and ECo,2 1c = 1.33 V)
that K(AgCl)>K(AgBr)>K(AgI)
(a)Er1Aagl1Ag
(a) 0.46 V
(b)0.60V
(c)0.53V
(d) -0.17 V
(b)E$2$1Ael1Ag
14.The EMF of the cell:Zn-Hg (C,M)|
Zn$2$(aq)| Zn-Hg (CM) at 25$/circ$C, if the
(d)E IAel1Ag=
=E$/circ$
concentrations of the zinc amalgams are:
C,= 10 g per 100 g of mercury, C2= 1 g
10.The EMF of a galvanic cell composed
per 100 g mercury,is
of two hydrogen electrodes is 272 mV.
(a) 0.059 V
(b) 0.0295 V
What is the pH of the solution in which
(c) 0.59 V
(d) 0.295 V
the anode is immersed if the cathode is in
contact with a solution of pH=3?
\end{problembox}

\begin{problembox}
\textbf{Problem 33} \hfill \textbf{[Pearson Advanced Electrochemistry Q15]}
\label{q:ext-pearson-33}

What is the equilibrium constant of the
reaction: 2Fe$2$++Au*$/to$2Fe$2$++Au$2$?
(a) 3
(b) 6.7
Given EaIA = 1.68 V, Eas*IA = 1.50 V,
(c) 7.6
(d) 1.6
EpIFa/\# = 0.75 V and 2.303RT/F = 0.06.
11.At what pH does the potential (EMF)
for the disproportionation of chlorine
(a) 1$/times$1022
changefrom a negative value to a positive
(b) 1 $/times$ 10-22
value, assuming 1.0 bar pressure and
(c) 1$/times$ 10-1
1.0 M concentration for all species except
(d) 1 $/times$10-72
hydrogen ion?Given:
\end{problembox}

\begin{problembox}
\textbf{Problem 34} \hfill \textbf{[Pearson Advanced Electrochemistry Q16]}
\label{q:ext-pearson-34}

What is the EMF of the cell:Pt,H2
Cl+2e$/to$2CI;E$/circ$=1.36V
(1 atm)|CH,COOH (0.1 M) II (0.01
++H+O
M)NHOH H,(1 atm),Pt?Given:
E$/circ$= 1.63 V.
Kfor CH,COOH =1.8 x 10-,K, for
NHOH = 1.8 $/times$ 10-, 2.303RT/F = 0.06,
[2.303RT/F = 0.06]
1og 1.8 = 0.25)
(a) 4.5
(b) 1.5
(a) 0.465 V
(b)-0.465 V
(c) 2.25
(d) 9.0
(c)-0.2325V
(d)-0.93V
--- PAGE 16 ---
8.16Chapter8
17.Two electrochemical cells are assembled
-0.018 V,respectively,at 25$/circ$ C.If the
in which the following reactions occur:
potential of normal calomel electrode is
V2++VO$2$+ 2H+$/to$2V+HO;Ecl
-0.28 V, the EMF of the Daniel cellis
= 0.616 V
(a) 1.065 V
(b) 1.101 V
V3++Ag*+H,O $/to$VO$2$++Ag(s)+ 2H";
(c)0.803V
(d) 0.262V
E$/circ$ce=0.439V
22.A hydrogen electrode placed in a buffer
If E'g =0.799 V, what is E?
solutionof CH,COONa and CH,COOH
in the ratio's x:y and y:x has electrode
(a)-0.256V
(b)+0.256 V
potential values E, and E, volts,
(c) +1.854 V
(d) -1.854 V
respectively,at 25$/circ$C.The pK. value of
18.The EMF of cell: H(g)|BufferlNormal
acetic acid is
calomelelectrode,is 0.70Vat 25$/circ$C,when
(a) (E,+ E)/ 0.118
barometric pressure is 760 mm.What is
(b)(E-E)/0.118
the pH of the buffer solution?E$/circ$Calomel
(c)-(E,+E)/0.118
= 0.28 V.[2.303RT/F = 0.06]
(d)(E,-E)/0.118
(a) 3.5
23.The standard reduction potential at 25$/circ$C
(b) 7.0
of the reaction 2HO+2e=H+2OH
(c) tending to zero
is-0.84V.Calculate equilibrium constant
(d) tending to 14.0
for thereaction:2HO=HO++OHat
25$/circ$C.(2.303RT/F=0.06)
19.Whatisthesolubilityproductofa
saturated solution of AgCrO, in water
(a) 10-14
(b)1014
at 298 K if the EMF of the cell: Ag|Ag*
(c)102
(d) 10-2
(satd. AgCrO)I Ag* (0.1 M)| Ag is 0.162
\end{problembox}

\begin{problembox}
\textbf{Problem 35} \hfill \textbf{[Pearson Advanced Electrochemistry Q24]}
\label{q:ext-pearson-35}

The voltage of the cell given below is
V at 298 K?[2.303RT)F = 0.06,1og 2 = 0.3]
0.61 V.
(a) 2.0 $/times$ 10-4
(b) 3.2$/times$10-i
(c)8.0$/times$10-12
(d) 4.0$/times$10-12
Pt|H(1 bar)|NaHSO, (0.4 M), Na,SO;
(6.4$/times$ 10-2 M) IIZn$2$+ (0.4 M)|Zn
20.A TI|Tlcouple was prepared by saturating
If E$/circ$z$2$zn=-0.76 V, Calculate Ka2 of
0.10 M -KBr with TIBr and allowing
H,SO(2.303RT/F=0.06)
Tl ions form the insoluble bromide to
equilibrate.This couple was observed to
(a) 3.2 $/times$ 10-4
(b) 3.2$/times$10-2
have a potential-0.444 Vwith respect
(c) 3.2 $/times$ 10-3
(d) 6.4$/times$ 10-7
to Pb$2$t|Pb couple in which Pb$2$+ was
25.B
Estimate thecell potentialof aDanielcell
0.10 M. What is the K, of TIBr.[Given:
having 1.0 M-Zn$2$+ and originally having
Ep*Po = -0.126 V, ErrIm = -0.336 V,
1.0 M-Cu$2$+after sufficient ammonia has
log 2.5 = 0.4, 2.303 RT/F = 0.06]
been addedtothe cathode compartment
(a) 4.0$/times$10-6
(b) 2.5$/times$10-4
to make the NH, concentration 2.0 M.
Given: E$/circ$z2/zn = -0.76 V, E$/circ$cu$2$cu
(c) 4.0$/times$10-s
(d) 6.3$/times$10-3
= +0.34 V, Kfor Cu (NH)+ = 1 $/times$ 1012.
21.The EMFs of the cell obtained by
(a) 1.10 V
combining separately Zn and Cu
(b) 0.704 V
electrodes of a Daniel cell with normal
(c) 0.396 V
calomelelectrodesare1.083Vand
(d) 1.496 V
--- PAGE 17 ---
Electrochemistry8.17
26.The minimum mass of NaOH required
measured between the electrodes at 25$/circ$C?
tobe added in RHS to consume all the
H' present in RHS of the cell of EMF,
= 10-$/circ$)
+0.70 V at 25$/circ$C, before its use.
(a) 0.92 V
(b) 0.32 V
Zn|Zn$2$+(0.01 M)IIHC1(V=500 ml)|
(c) 1.28 V
(d) 0.56 V
H(1 bar)|Pt
31.Although aluminium is above hydrogen
Given: E$/circ$zn$2$tzn = -0.76 V; 2.303RT/F
inthe electrochemical series,it is stablein
= 0.06
air and water. Why?
(a) 2.0 g
(b) 4.0 g
(a)Aluminium is non-reactive metal.
(c) 0.2 g
(d) 0.4 g
(b)A layer of oxide is formed at metal
surface,which preventfurtherreaction.
27.The EMF of the cell:Ag,AgCl in 0.1 M
- KCIII satd. NHNO,I0.1 M - AgNO,
(c)The reaction is thermodynamically
unfavourable,i.e.,Gis not negative.
Ag is 0.42 V at 25$/circ$C. 0.1 M -KClis 50/%
(d)Kinetically, the reaction has very low
dissociated and 0.1 M -AgNO, is 40/%
dissociated. The solubility product of
activation energy andhence do not
AgCIis (2.303RT/F =0.06)
occur.
(a) 1.0$/times$10-10
32.When silver chloride is dissolved in a
(b) 2.0 $/times$ 10-9
large excess of ammonia,practically all
(c)1.0$/times$10-9
silver ion can be assumed to exist in form
(d) 2.0$/times$10-10
of a single ionic species [Ag(NH),]
Compute the values of x and y using the
\end{problembox}

\begin{problembox}
\textbf{Problem 36} \hfill \textbf{[Pearson Advanced Electrochemistry Q28]}
\label{q:ext-pearson-36}

On the basis of information available
twofollowing cells:
from thereaction
Cell I: Ag | 4.0 $/times$ 10-4 M - AgCl, 1 M
4A1+3O2$/to$2AlO;$/Delta$G=-965kJ/mol
- NH, II 4 $/times$ 10-2 M - AgC1, 1 M
of O2
-NH|Ag; Een = 0.118 V at 298 K.
The minimum EMF required to carry out
CellII:Ag |3 $/times$ 10-$2$ M -AgCl, 1 M
electrolysis of AlOis
- NH, II 3 $/times$ 10-* M - AgC1, 0.1 M
(a) 0.833 V
(b) 2.5 V
-NH|Ag; Ecn = 0.118 V at 298 K.
(c) 5.0 V
(d) 1.67 V
(a) x=1,y=2
(b) x= 1,y=4
29.The theoreticalefficiency of ahypothetical
(c)x=2,y=4
cellisabout84/%whichinvolvesthe
(d) x=1,y=6
followingreaction
A(s)+ B$2$(aq)
)$/to$A$2$(aq)+B(s);
33.When metallic copper is shaken with a
$/Delta$H = -285 kJ,
solution of a copper salt, the reaction
Cu+Cu$2$+2Cu*proceeds.When
then,the standard EMF of the cell is
equilibrium is established at 298 K,
(a) 1.10 v
(b)1.24 V
[Cu$2$/[Cu']= 1.667 $/times$ 10$/circ$M-'.If the
(c) 2.48 V
(d) 2.20 V
standard potential of the Cu$2$|Cu half-
\end{problembox}

\begin{problembox}
\textbf{Problem 37} \hfill \textbf{[Pearson Advanced Electrochemistry Q30]}
\label{q:ext-pearson-37}

To a beaker containing 0.1 M - HCl, a
cellis +0.3376V,what is the standard
little pure solid AgCl is added. Both
potential of Cu| Cu half-cell?(Given:
a silver and a hydrogen electrode (PH,
2.303RT/F=0.06,1og 2=0.3,log 3 =0.48)
=1.0 bar) are then placed in the solution.
(a)-0.3732 V
(b) 0.6752V
What is the approximate value of EMF
(c)0.5242V
(d) 0.151 V
--- PAGE 18 ---
8.18$/cdot$Chapter8
\end{problembox}

\begin{problembox}
\textbf{Problem 38} \hfill \textbf{[Pearson Advanced Electrochemistry Q34]}
\label{q:ext-pearson-38}

Zinc granules are added in excess to a
Use thefollowingdatatocalculate$/Delta$G$/circ$
500 ml of 1.0 M nickel nitrate solution
for the above reaction.
at 25$/circ$C until the equilibrium is reached.
K[Au(CN)]
If the standardreduction potential of
= X
Zn$2$Zn and Ni$2$|Ni are -0.75 Vand
O2+2HO +4e $/to$4OH'; E$/circ$=+0.41 V
-0.24 V, respectively, the concentration
Au$2$+ 3e"$/to$Au;
E$/circ$=+1.50 V
of Ni$2$in solution at equilibrium is
Au$2$++ 2e $/to$ Au*;
E$/circ$=+1.40 V
(2.303RT/F= 0.06)
(a)-RTInX+ 1.29 F
(a) 1.0 $/times$10-" M
(b)-RTlnX-1.29 F
(b) 1.0$/times$10" M
(c)+RTInX+ 2.11 F
(c)5$/times$10-M
(d)-RTInX-2.11 F
(d) 2$/times$10-1 M
35.The EMF for the cell: Ag(s)|AgCl(s)|
\end{problembox}

\begin{problembox}
\textbf{Problem 39} \hfill \textbf{[Pearson Advanced Electrochemistry Q39]}
\label{q:ext-pearson-39}

.F2 gas cannot be obtained by the
KCI(0.2 M)II KBr (0.001 M)|AgBr(s) |
electrolysis of any aqueous fluoride salt
Ag(s) at 25$/circ$C is (K(AgCl) = 2.0 $/times$ 10-10;
because
K(AgBr)=4.0$/times$10$/circ$i3,2.303RT/F=0.06,
(a)F is the strongest oxidizing agent
log 2 = 0.3)
(b)Feasily combines with water
(a) 0.024 V
(c)F, readilycombineswiththe
(b) -0.024 V
electrodes
(c)-0.24 V
(d)F" can never be oxidized
(d) -0.012 V
40.A volume of 100 mlof a buffer of 1 M
36.Two students use same stock solution of
NH, and 1 M -NH*is placed in two
ZnSObut different solutions of CuSO
half-cells connected by a salt bridge.A
The EMF of one cell is 0.03 V higher
current of 1.5 A is passed through the cell
thanthe other.The concentrationof
for 20 min.If electrolysis of water takes
CuSOin the cell with higher EMF value
place only and the electrode reactions are:
is 0.5 M.The concentration of CuSOin
the other cell is (2.303RT/F = 0.06)
Right: 2H,O +O+4e $/to$4 OH
(a) 0.05 M
(b) 5.0 M
and Left: 2HO$/to$4H*+O2+4e,
(c) 0.5 M
(d) 0.005 M
then, the pH of the
37.Two weak acid solutions HA,and HA
(a)right electrode will increase
eachwiththe same concentration and
(b) left electrode will increase
having pK.  values 3 and 5 are placed in
(c)both electrode will increase
contact with hydrogen electrodes (1 atm,
(d) both electrode will decrease
25$/circ$C)and are interconnected through a
41.When molten ICI, is electrolysed using
salt bridge.EMF of the cellis
platinum electrodes
(a) 0.0295 V
(b) 0.118 V
(c)0.0885V
(d)0.059V
(a) I, is evolved at cathode and Cl, at
anode
38.Consider the reaction of extraction of
(b) Cl, is evolved at cathode and I at
gold from its ore:
anode
Au(s) + 2CN (aq)+
1
(c) I is evolved at cathode and both I2
O(g)+
4
and Cl,at anode
HO(I)$/to$Au(CN) (aq)+OH(aq)
(d)electrolysis does not take place
--- PAGE 19 ---
Electrochemistry8.19
42.When electric current is passed through
47.Which of the following substances: Na,
a cell having an electrolyte,positive ions
Hg, S, Pt and graphite can be used as
move towards the cathode and negative
electrodes in electrolytic cells having
ions towards the anode.If thecathodeis
aqueous solution?
pulled out of the solution,then the
(a)Na, Pt and graphite
(a) positive and negative ions will move
(b)Na and Hg
towards the anode
(c)Hg, Pt and graphite
(b)positive ions will
start
tmoving
(d) Na and S
towards the anode, the negative ions
\end{problembox}

\begin{problembox}
\textbf{Problem 40} \hfill \textbf{[Pearson Advanced Electrochemistry Q48]}
\label{q:ext-pearson-40}

3.An ammeter anda copper voltameter
will stop moving
are connected in series in an electric
(c)negative ions will continue to move
towards the anode, the positive ions
circuit through which a constant direct
current flows.The ammeter shows 0.5 A.
will start moving randomly
If 0.635 g of Cu is deposited in 0.965 h,
(d) positive and negative ions will start
what is the percentage error of ammeter?
moving randomly
(Cu = 63.5)
\end{problembox}

\begin{problembox}
\textbf{Problem 41} \hfill \textbf{[Pearson Advanced Electrochemistry Q43]}
\label{q:ext-pearson-41}

When 10-6 M -HClis electrolysed
(a) 5 /%
(b) 10 /%
(a) Ois produced at the anode
/%6 ()
/%06 (P)
(b)H,is produced at the anode
49.The same current is passed through
(c) Cl, is produced at the anode
acidulatedwater andstannous chloride
(d) Cl,and Oare produced at the anode
solution.What volume of dry detonating
gas at 0$/circ$C and 1 atm is evolved from
44.A dilute aqueous solutionof Na,SOis
electrolysed using platinum electrodes.
water, when 1.20 g of tin is deposited
The products at the anode and cathode
from the other solution?(Sn=120)
are,respectively,
(a) 112ml
(b)336ml
(a) O2, H2
(c)224 ml
(d) 672 ml
(b) S,O$2$-, Na
50.The preparation of LiOH by the
(c) O2, Na
electrolysis of a 35/% solution of LiCl
(d) S,O$2$,H
using a platinum anode led to a current
efficiency of 80/%.What weight ofLiOH
45.In the electrolysis of a fused salt, the mass
was formed by the passage of 2.5 A for
of substance deposited on an electrode
4825 s?
will not depend on
(a) 1.92 g
(b) 2.40 g
(a)temperature of bath
(c) 0.96 g
(d) 0.672 g
(b) current intensity
(c) time of electrolysis
\end{problembox}

\begin{problembox}
\textbf{Problem 42} \hfill \textbf{[Pearson Advanced Electrochemistry Q51]}
\label{q:ext-pearson-42}

Assuming that copper contains only
(d)electrochemical equivalent of ions
iron, silver and gold as impurities.
After passage of 12.4 A for 4825 s, the
46.When an electric current is passed through
mass of anode decreased by 20.00 g
an aqueous solution of the following,the
and the cathode increased by 19.05 g.
concentrations of cation as well as anion
The percentages of iron and copper in
will not change for (assume constant
the original sample are, respectively,
volume of the electrolytic solution)
(Cu=63.5,Fe=56)
(a) CsC1
(a)4.75/%,95.25/%
(b)KNO,
(b)2.8/%,95.25/%
(c) AgNO;
(c)95.25/%,4.75/%
(d)HCI
(d) 95.25/%,1.91/%
--- PAGE 20 ---
8.20Chapter8
\end{problembox}

\begin{problembox}
\textbf{Problem 43} \hfill \textbf{[Pearson Advanced Electrochemistry Q52]}
\label{q:ext-pearson-43}

After electrolysis of an aqueous sodium
57.Most of the copper used to make wire has
chloride solution,it was found that the
been electrically refined by depositing it
solution is being neutralized by 60 ml N
from copper salts solution (divalent) on
-HCl solution.During the same period
to a cathode.What is the cost of electrical
of electrolysis, 3.18 g of copper was
energy required per kg of copper if the
deposited in a copper voltameter in series.
cost of electricity is Rs 4.00 per kWh
What is thepercentage of the theoretical
and the cell operates at 0.33 V? The
yield ofsodiumhydroxideobtained?
electrochemical equivalent of copper is
(Cu=63.6)
0.00033 g/coulomb.
(a) 40/%
(b) 60/%
(a) Rs 11.11
(b)Rs 5.55
(c)30/%
(d)80/%
(c) Rs 2.22
(d) Rs 1.11
53.A lead storage batteryhas initially 200g
\end{problembox}

\begin{problembox}
\textbf{Problem 44} \hfill \textbf{[Pearson Advanced Electrochemistry Q58]}
\label{q:ext-pearson-44}

Electrolysisofan
1acetatesolution
of lead and 200 g of PbO2 plus excess
produces
ethaneaccordingtothe
H,SO.Theoretically,how long could this
reaction:
cell deliver a current of 10.0 A,without
2CHCOO-$/to$CH(g)+2CO(g)+2e
recharging,if it were possible to operate
it so that the reactiongoes to completion?
Whattotalvolume of ethane andCO2
(Pb =208)
would be produced at 0$/circ$C and 1 atm,if
a current of 0.5 A is passed through the
(a) 16083.33 s
(b) 8041.67 s
solution for 482.5 min?Assume current
(c) 44.68 s
(d) 18557.7 s
efficiency 80/%.
54.A volume of 100 ml of 0.6 N-CuSO4
(a) 1.344 L
(b) 2.688 L
solutioniselectrolysedbetweentwo
(c) 4.032 L
(d) 1.792 L
platinum electrodes till the concentration
\end{problembox}

\begin{problembox}
\textbf{Problem 45} \hfill \textbf{[Pearson Advanced Electrochemistry Q59]}
\label{q:ext-pearson-45}

To perform an analysis of a mixture of
in the residual liquid is 0.1 N,when a
steady current of 5.0 A is used. How long
metalions by electro-deposition,the
second metal to be deposited must not
should the current be passed to get the
above change?
being plating out until the concentration
ratio of the second to thefirstis about 10$/circ$.
(a)965 s
(b)96500 s
What must be the minimum difference
(c) 1930 s
(d) 482.5 s
in standard potential of the twometals
55.A constant current flowed for 2h through
which form dipositive ions in order for
a potassium iodide solution, oxidizing
such an analysis to be feasible?
theiodideiontoiodine.Attheendofthe
(a) 0.177 V
(b) 0.354 V
experiment,the iodine was titrated with
(c) 0.708V
(d) 0.088V
72 ml of 1.0 M-Na,SO, solution.What
\end{problembox}

\begin{problembox}
\textbf{Problem 46} \hfill \textbf{[Pearson Advanced Electrochemistry Q60]}
\label{q:ext-pearson-46}

During the electrolysis of 0.1 M-CuSO
was the averagerate ofcurrent flow,in
solution using copper electrodes,a
ampere?
depletion of Cu$2$occurs near the cathode
(a) 0.965 A
(b)1.93 A
with a corresponding excess near the
(c) 0.483 A
(d) 0.0965 A
anode, owing to inefficient stirring of
56.Anthracene, C,4Hio, can be oxidized to
the solution.If the local concentration
anthraquinone,CHO.What weight of
of Cu$2$+ near the anode and cathode are,
anthraquinone can be produced by the
respectively, 0.12 M and 0.08 M, the
passage of a current of 1 A of 40 min if
back EMF developed at 298 K is (log 1.5
the current efficiency is 96.5/%?
=0.18,2.303RT/F= 0.06)
(a) 0.862 g
(b) 0.832 g
(a) 0.33 V
(b) 5.4mV
(c) 0.624 g
(d) 0.738 g
(c) 2.7mV
(d)10.8mV
--- PAGE 21 ---
Electrochemistry8.21
61.The following galvanic cell:
The first reactionis anelectrolytic
Zn|Zn(NO)(100 ml, 1 M)Cu(NO3)
reaction and second is steam distillation.
What amount of current would have tobe
(100 ml,1 M)| Cu
used infirst reaction to produce enough
was operated as an electrolytic cell as Cu
intermediate to yield102g pure HOper
as the anode and Zn as the cathode.A
hour.Assume current efficiency50/%.
current of 0.4825 A was passed for 10 h
(a) 643.33 A
and thenthecellwas allowedtofunction
(b) 321.67 A
as galvanic cell. What would be the final
(c) 160.83 A
EMF of the cell at 25$/circ$C? Assume that
the only electrode reactions occurring
(d)1286.67 A
were those involving Cu|Cu$2$and Zn|
Zn$2$. Given: E$/circ$$2$cu = +0.34 V, Ez$2$zn
65.In an analytical determination of arsenic,
=-0.76 V, 2.303RT/F= 0.06,1og 1.9= 0.28)
asolution containing arsenious acid,
(a) 1.10 v
(b) 1.0616 V
HAsO,KIand a smallamountofstarch
(c) 1.1384 V
(d) 1.1768 V
is electrolysed.The electrolysis produces
free IfromIions and the Iimmediately
\end{problembox}

\begin{problembox}
\textbf{Problem 47} \hfill \textbf{[Pearson Advanced Electrochemistry Q62]}
\label{q:ext-pearson-47}

At the Nangal fertilizer plant in Punjab,
oxidizes the arsenious acid to hydrogen
hydrogen is produced by the electrolysis
arsenate ion, HAsO$2$.
of water. The hydrogen is used for the
production of ammonia and nitric acid
I+HAsO+HO$/to$ 2I+HAsO$2$
(by the oxidation of ammonia).If the
+ 4H*
averageproductionof ammonium nitrate
is 1152 kg/day, the daily consumption of
Whentheoxidationofarsenicis
complete, the free iodine combines with
electricity (in A /day)is
the starch to give a deep blue colour. If
(a)96,500
(b) 48,250
during a particular run, it takes 96.5 s
(c) 32,166.67
(d)24,125
for a current of 10.0 mA to give an end
\end{problembox}

\begin{problembox}
\textbf{Problem 48} \hfill \textbf{[Pearson Advanced Electrochemistry Q63]}
\label{q:ext-pearson-48}

Lactic acid, HCH,O, produced in 1 g
point (indicated by the blue colour),
sample of muscle tissue was titrated
how many g of arsenic are present in the
using phenolphthaleinas indicator
solution.(As =75)
against OHions whichwere obtained by
(a) 6.3 $/times$10-4g
theelectrolysis ofwater.As soon as OH
(b) 1.26$/times$ 10-g
ions are produced, they react with lactic
(c) 3.15$/times$ 10-g
acidandatcompleteneutralization,
(d) 6.3$/times$10-g
immediatelyapink colour is noticed.
If electrolysis was made for 1158 s using
50.0 mA current to reach the end point,
\end{problembox}

\begin{problembox}
\textbf{Problem 49} \hfill \textbf{[Pearson Advanced Electrochemistry Q66]}
\label{q:ext-pearson-49}

ElectrolysisofasolutionofMnSO
in aqueous sulphuric acid is a method
what was the percentage of lactic acid in
for the preparation of MnO as per the
muscle tissue?
reaction:Mn$2$(aq)+2HO $/to$MnO2(s)
(a) 5.4/%
(b) 2.7/%
+ 2H(aq) + H(g). Passing a current of
(c) 10.8/%
(d) 0.054/%
19.3 A for 2 h gives only 52.2 g of MnO2.
64.HO2 can be prepared by successive
The current efficiency is (Mn=55)
reactions:
(a)8.33/%
2NHHSO$/to$H,+(NH)SO
(b) 83.33/%
(NH)SO+ 2HO $/to$ 2NHHSO4
(c) 41.67/%
+H,O2
(d)100/%
--- PAGE 22 ---
8.22Chapter8
67.Copper sulphate solution (250 ml) was
Pb +SO$2$$/to$PbSO4+ 2e
electrolysed using a platinum anode and
PbO2+ 4H*+ SO$2$-+ 2e $/to$PbSO4
a copper cathode. A constant current
+ 2HO
of 2 mA was passed for 19.3 min. It
was found that after electrolysis the
(a) 965
(b) 482.5
absorbance of the solution was reduced
(c)1930
(d) 241.25
to 50/%of its original value.Calculate the
71.A cell whose resistance, when filled with
concentration of copper sulphate in the
solution to begin with.
0.1 M-KCl is 200 Q, is measured to be
6400 Q, when filled with 0.003 M -NaCl
(a) 9.6$/times$10- M
solution.What is the molar conductance
(b) 4.8$/times$10- M
of NaCl solution, in Q-'cm$2$mol-' if the
(c) 2.4$/times$ 10-s M
molar conductance of 0.1 M -KCl is
(d) 1.2 $/times$ 10- M
120 Q-'cm$2$eq-?
68.(
Calculate the quantityofelectricity that
(a) 41.67
(b) 250
would be required to reduce 12.3 g of
(c)125
(d)375
nitrobenzene to aniline if the current
efficiency for the process is 50/%. If the
72.At 18$/circ$C,the ionic mobilities of NH4and
potential drop across the cellis 3.0V,how
ClOions are 6.6x 10- and 5.4 $/times$ 10-8
much energy will be consumed?
m$2$volt-'s-' at infinite dilution.What is
(a) 0.3 F,86.85 kJ
themolar conductanceofammonium
chromate solution inQ-'m$2$mol-?
(b) 0.6 F, 173.7 kJ
(c)1.2 F, 347.4 kJ
(a) 0.01795
(b) 0.01158
(d) 0.6 F,173.7 kJ
(c) 1.2$/times$ 10-8
(d) 1.86 $/times$ 10-7
69.Perdisulphuric acid, H,SOg can be
\end{problembox}

\begin{problembox}
\textbf{Problem 50} \hfill \textbf{[Pearson Advanced Electrochemistry Q73]}
\label{q:ext-pearson-50}

At 25$/circ$C,the molar conductance at
preparedby electrolytic oxidation of
infinitedilutionfor HClsolutionis
H,SO, oxygen and hydrogen gases are by
4.25 $/times$ 10-$2$Q-' m$2$mol', while its specific
products. In such an electrolysis, 4.54 L
conductance is 382.5Q-'m-'.If the degree
of H, and 2.27 L of Owere generated at
of dissociation is 90/%,the molarity of
STP. What is the mass of H,SOformed?
solution is
(a) 0
(a) 0.9 M
(b) 77.6 g
(b) 1.0 M
(c) 38.8 g
(c) 0.1 M
(d) 19.4 g
(d) 1.1 M
\end{problembox}

\begin{problembox}
\textbf{Problem 51} \hfill \textbf{[Pearson Advanced Electrochemistry Q70]}
\label{q:ext-pearson-51}

During the discharge of a lead storage
74.A big irregular shaped vessel contained
battery,the density of sulphuric acid fell
water,conductivity of which was 2.56
from 1.5 to 1.1 g/ml. Sulphuric acid of
$/times$ 10-3 s-'m-'.585 g of NaCl was then
density 1.5 g/ml is 40/% HSO4,by weight,
addedtothewater andconductivityafter
and that of density 1.1 g/mlis 10/% HSO4,
the additionof NaCl,wasfound tobe3.06
by weight.The battery holds 3.6 L of the
$/times$10-3 s-'m-'.The molar conductivity of
acid and the volume remained practically
NaCl at this concentration is 1.5 x 10-2
constant during the discharge.Calculate
S-'m$2$mol-'. The capacity of vessel if it is
thenumber ofampere-hourswhich
fulfilled with water,is
the battery should have been used.The
(a)3$/times$1041
(b) 301
electrodereactions are:
(c) 3$/times$1081
(d) 3$/times$10$/circ$1
--- PAGE 23 ---
Electrochemistry8.23
75.Themolarconductivityof0.10M
78.B
Resistanceof0.2Msolutionof
solution of MgCl is 100 mho cm$2$mol-',
an electrolyte is 50 Q.The specific
at 25$/circ$C.A cell with electrodes that are
conductance of the solution is 1.4 S m-.
1.50 cm$2$ in surface area and 0.50 cm
The resistance of 0.5 M solution of the
apart is filled with 0.10 M - MgCl
same electrolyte is 280Q.The molar
solution.How much current will flow
conductivity of 0.5M solution of the
when the potential differencebetween the
electrolyte (in S m$2$mol-')is
electrodesis5volts?
(a) 5$/times$10-3
(b) 5$/times$103
(a) 0.03 A
(b) 3.0 A
(c) 5$/times$102
(d) 5$/times$10-4
(c) 0.15 A
(d) 15 A
\end{problembox}

\begin{problembox}
\textbf{Problem 52} \hfill \textbf{[Pearson Advanced Electrochemistry Q79]}
\label{q:ext-pearson-52}

The conductivity ofa saturated solution
76.For Na*, the value ofsymbolis
containing AgA(K=3$/times$10-14)and AgB
50.0 Q-' cmmol-'.The speed of Naion
(Ks = 1 $/times$ 10-14) is 3.75 $/times$ 10-8Q-cm-.
in the solution,if in the cell,electrodes
If the limiting molar conductivity of
are 5 cm apart and to which a potential
Ag* and A ion is 60 and 80 Q-icm$2$
of 19.3 volt is applied is
mol',respectively,the limiting molar
(a) 2$/times$10-3cm/s
conductivity of B(in Q-'cm$2$mol-')is
(b) 1 $/times$10-3cm/s
(a)135
(c) 2$/times$10-4cm/s
(b) 67.5
(d) 2$/times$10-$2$cm/s
(c)270
77.The conductivity of a saturated solution
(d) 190
of AgCl at 298 K was found to be 3.40
$/times$10-6Q- cm-'; the conductivity of water
80.The conductivity of saturated solution
used to make up the solution was 1.60
of Ba(PO) is 1.2 $/times$ 10- Q-  cm-. The
$/times$10-6Q- cm-'.Determine the solubility
limiting equivalent conductivities of
of AgClin water in mole per litre at 298 K.
BaCl,KPO4and KCl are 160,140
and i00 Q-  cm$2$ eq- , respectively. The
The equivalent conductivity of AgCl at
infinite dilution is 150.0Q- cm-$2$eq-
solubility product of Ba,(PO) is
(a) 1.44 $/times$ 10-10
(a) 10-s
(b) 1.2 $/times$ 10-5
(b) 1.08$/times$10-23
(c) 3.33$/times$10-s
(c)1.08$/times$10-25
(d) 1.2 $/times$ 10-8
(d) 1.08 $/times$ 10-27
SectionB(OneorMorethanoneCorrect)
1.Whichof thefollowingstatement(s)
2.Pick up the false statement(s):
differentiatebetween electrochemicalcell
(a)The net chemicalchange in a galvanic
andelectrolytic cell?
cell reaction is always redox reactions.
(a) Spontaneous or non-spontaneous
(b)In a galvanic cell made of cobalt and
nature of the chemical process.
cadmium electrodes,cobalt electrode
(b) Chemical reactions occurring at the
acts as anode.
electrodes.
(c) Standard potential increases with
(c)Positive and negative nature of
increasing
concentrationofthe
anode.
electrolyte.
(d)Dependence on Faraday's law.
(d) Calomel electrode is a reference
electrode having 0.00 volt potential.
--- PAGE 24 ---
8.24Chapter8
3.Consider the cell: Ag(s),AgCl(s)|KCl
8.The function(s) of salt bridge in a cell
(0.1 M) | HgzCl,(s), Hg(l). The cell
is/are
potential
(a) It maintains standard electrode
(a)increases on increasing concentration
potential of cell constant which
of Cl ions.
depends on severalfactors.
(b)decreases on decreasingconcentration
(b)It completes the electrical circuit.
of CI ions.
(c)It departs both the solutions from
(c)is independent of concentration of
each other.
CI ions.
(d)It maintains the electrical neutrality
(d) is independent of amounts of AgCl
of both electrolytic solutions.
and HgCl2.
9.The reactions taking place in the dry cell
4.The passage of electricity in the Daniel
are:
cellwhenZnandCuelectrodesare
Anode:Zn$/to$Zn$2$++2e
connected
Cathode:2MnO+2NH4*+2e
(a)from Cu to Zn inside the cell
$/to$MnO + 2NH, + HO
(b)from Cu to Zn outside the cell
(c)from Zn to Cu inside the cell
The minimum mass of reactants, if a dry
(d)from Zn to Cu outside the cell
cell is to generate 0.25A for 9.65 h, are
(Mn = 55,Zn = 65.4) (neglect any other
5.From an electrolyte,one mole of electron
chemical reactions occurring in the cell)
will deposit at cathode
(a) 2.943 g Zn
(b)7.83g MnO2
(a) 63.5gof Cu
(b) 24 g of Mg
(c)1.62 g NH4*
(d) 3.915g MnO2
(c) 11.5 g of Na
(d) 9.0 g of Al
\end{problembox}

\begin{problembox}
\textbf{Problem 53} \hfill \textbf{[Pearson Advanced Electrochemistry Q10]}
\label{q:ext-pearson-53}

EMF of the cell:Cd(s)|CdCl.5HO
6.Which of the following statements is/are
(sat.)|AgCl(s)|Ag(s) is +0.70 V at 0$/circ$C
correct?
and +0.60 V at50$/circ$C.If $/Delta$H$/circ$and $/Delta$S$/circ$
(a) The conductivity of molten NaCl is
are temperature independent,then the
due to movement of Na+ and Cl ions.
correct information(s) regarding the cell
(b) Solid NaCl is good conductor of
reactionis/are
electricity.
(a) $/Delta$G$/circ$= -115.8 kJ at 50$/circ$C
(c)Molten sodium is a good conductor
(b) $/Delta$G$/circ$= 135.1 kJ at 0$/circ$$/circ/text(C)$
because of mobile electrons.
(c) $/Delta$S$/circ$=-386J/K
(d)Resistivity is reciprocal of molar
(d)$/Delta$H$/circ$= -221.178 kJ
conductivity of electrolyte.
11.Two litre solution of a buffer mixture
7.Ionic conductance at infinite dilution of
containing 1.0 M-NaHPO and 1.0 M
A13 and SO$2$- ions are 60 and 80 Q-1
-NaHPOis placed in two compartments
cm$2$eq',respectively.The correct detail(s)
(one litre in each) of an electrolytic cell.
regardingAl(SO)is/are
The platinum electrodes are inserted in
(a) the molar conductance is140
each compartment and 1.25 A current is
Q-'cm$2$mol'.
passed for 965 min.Assuming electrolysis
(b) the equivalent conductance is 140
ofonlywater ateachcompartment,
Q-' cm$2$ eq-'.
what will be pHin each compartment
(c)the molar
conductance
: is840
after passage of above charge?(pK. for
Q-'cm$2$mol-'.
HPO= 2.15,1og 7=0.85)
(d) the molar conductance is 23.33
(a)Anode:3.00
(b)Cathode:3.00
Q-'cm$2$mol-'.
(c)Anode:1.30
(d) Cathode:1.30
--- PAGE 25 ---
Electrochemistry8.25
12.For the reduction of NO, ion in an
14.In a galvanic cell, the salt bridge
aqueous solution, E$/circ$ is +0.96 V. Values of
(a) does not participate chemically in the
E$/circ$for some metalions are given below:
cell reaction.
v2t(aq)+2e$2$$/to$V(s)
E$/circ$=-1.19 V
(b) stops the diffusion of ions from one
electrodeto another.
Fe$2$(aq)+ 3e"$/to$ Fe(s)
E$/circ$=-0.04V
(c)isnecessaryfor the occurrenceof the
cell reaction.
Au$2$(aq)+ 3e$/to$Au(s)E$/circ$=+1.40 V
(d)ensuresmixingof thetwoelectrolytic
Hg$2$(aq)+2e$/to$Hg(l)E$/circ$=+0.86 V
solutions.
\end{problembox}

\begin{problembox}
\textbf{Problem 54} \hfill \textbf{[Pearson Advanced Electrochemistry Q15]}
\label{q:ext-pearson-54}

A lead storage cell is discharged which
The pair(s) of metal that is(are) oxidized
causestheH,SO4electrolyte tochange
by NO,in aqueous solution is(are)
fromaconcentrationof40/%byweight
(a)V and Hg
(density = 1.260 g/ml) to 28/%,by weight.
(b) Hg and Fe
The original volume of electrolyte was
(c) Fe and Au
1 L.Identify the correct statement(s):
(d)Fe and V
(a) The overall cell reaction is: Pb(s)
()0997(b)09*H7+(s)09d+
13.Among the following, the intensive
+ 2H,O(1).
propertyis(properties are)
(b) A total of 2.0 moles of H,SO is
(a)Molar conductivity
reacted.
(c)The total charge released from anode
(b)Electromotive force
of the cellis 1.93x10/textasciicircum() coulomb.
(c)Resistance
(d)The mass of electrolytic solution has
(d)Heat capacity
decreased.
Section C(Comprehensions)
Comprehension1
Suppose that the S.H.E.was arbitrarily assigned a value of 1.00 V for 2H*(aq)+ 2e $/to$H,(g).What
would this do to the observed voltage under standard condition for each of the following:(Given:
EZa
\end{problembox}

\begin{problembox}
\textbf{Problem 55} \hfill \textbf{[Pearson Advanced Electrochemistry Q1]}
\label{q:ext-pearson-55}

3.Zn-Cu Cell
(a) +0.34 V
(b) -0.76 V
(a) +1.10 V
(c) +1.76 V
(d)-0.34 V
(b) -1.10 V
\end{problembox}

\begin{problembox}
\textbf{Problem 56} \hfill \textbf{[Pearson Advanced Electrochemistry Q2]}
\label{q:ext-pearson-56}

E$/circ/text(C)$(u
(c)0.0V
(d) Indeterminate
(a) -0.76 V
(b) +1.34 V
(c) +1.76 V
(d) +0.34 V
--- PAGE 26 ---
8.26Chapter8
ComprehensionII
The standard reduction potential of theAglAg electrode at 298K is 0.80V.The solubility product of
AgI is 6.4 $/times$ 10-1 at 298 K. (2.303RTIF = 0.06)
4.The potential of Ag'|Ag electrode in a
6.The potential of I (0.04 M)/AglIAg
saturated solution of AgI at 298K is
electrode at 198 K is
(a)-0.314 V
(b) +0.314 V
(a)-0.088 V
(c)-0.172V
(d) +0.172 V
(b) +0.088 V
(c)-0.172 V
5.The standard reduction potential of
(d) +0.172 V
I]Agl|Ag electrode at 298 K is
(a) -0.314 V
(b)+0.314 V
(c) -0.172 V
(d) +0.172 V
ComprehensionI
The Edison storage cell is represented as Fe(s)|FeO(s)|KOH(aq)|NiO(s)|Ni(s).The half-cell
reactions are
NiO(s)+HO(1)+2e2NiO(s)+2OH;E$/circ$=+0.40V
FeO(s)+HO(I)+2e=Fe(s)+2OH;E=-0.87V
7.What is the cell reaction?
9.How does cell EMF depend on increasing
(a) NiO;(s)+ Fe(s) $/to$ 2 NiO(s)+ FeO(s)
the concentration ofKOH?
(b) 2NiO(s) +FeO(s)$/to$ NiO(s)+ Fe(s)
(a) increases
(c)NiO(s) +FeO(s)$/to$2NiO(s) +Fe(s)
(b)decreases
+ O2
(c)remains unaffected
(d)None of these
(d)none of these
8.What is the standard cell EMF?
10.Whatis the maximumamountof
(a) 1.27 V
electrical energy that can be obtained
(b) 0.47V
from one mole ofNiO?
(c) -1.27 V
(a) 2.54 J
(b) 245.11 kJ
(d)-0.47 V
(c) 122.56 kJ
(d) 61.28 kJ
ComprehensionIV
Thecellpotentialfortheunbalancedchemicalreaction:
Hg(aq)+NO,(aq)+3H,O*(aq)$/to$2Hg$2$(aq)+HNO(aq)+4HO(l)
is measuredunder standard conditions in the electrochemical cellshown in the accompanying diagram.
E$/circ$=0.02V
dish A
dish B
anode
cathode
--- PAGE 27 ---
Electrochemistry8.27
11.In which dish, the solution is acidic?
\end{problembox}

\begin{problembox}
\textbf{Problem 57} \hfill \textbf{[Pearson Advanced Electrochemistry Q13]}
\label{q:ext-pearson-57}

How long willit take to produce 0.10 mole
(a) Dish A
(b)Dish B
of HNO2by this reaction if a current of
(c)Both
(d)None
10 A passes through the cell?
12.How many moles of electrons pass
(a)965 s
(b) 193s
throughthecircuit when0.60moleof Hg2+
(c)1930 s
and 0.30 mole of HNO2 are produced in
the cell that contains 0.50 mole of Hg2+
(d) 482.5 s
and 0.40 mole of NO, at the beginning
of the reaction?
(a) 0.30
(b) 0.60
(c) 0.15
(d) 1.20
ComprehensionV
A fuel cell is the device to convert the energy of a fuel into electrical energy without theuse of heat
engine,where thefuelisburnt directly.Suchconversionsare possiblebecause the combustionreactions
are essentially redox reactions andhighly exothermic as well as highly exergonic.Electrical energy
can be obtained indefinitely from a fuel cell as along as the outside supply of fuel is maintained.In
hydrogen-oxygenfuel cell,thefollowingreactionstakeplace:
Anode reaction:
2H(g)+4OH(aq)$/to$4HO()+4e
Cathode reaction:
O(g)+2HO(I)+4e$/to$4OH(aq)
Overall reaction:
2H(g)+O(g)$/to$2HO()
The overall reaction has a value of$/Delta$H$/circ$=-285.8kJ/mol and$/Delta$G$/circ$=-237.39kJ/mol at 25$/circ$C.
14.Whatis thestandardEMFof thecell?
17.What is the approximate value of $/Delta$S$/circ$ for
(a) 0.615 V
(b) 1.23 V
thefuel cell reaction at 25$/circ$C?
(c) 2.46 V
(d) 0.74 V
(a) -0.1624 JK-1
15.What volume of gaseous H, (at STP),
(b) -162.4 JK-1
(c) +162.4 JK-1
whencombinedwithexcessOinthefuel
(d) +0.1624 JK-1
cell,is needed to produce 23.739 kJ of
useful work under standard conditions?
18.The theoreticalefficiency of thefuelcellis
(a) 2.27 L
(b) 4.54 L
given by
(c) 1.13 L
(d) 2.24 L
(a) 83.06/%
(b)100/%
16.S
Suppose the concentration ofhydroxide
ion in the cell is doubled at 298 K. The
(c) 67.53/%
(d) 97.88/%
cellvoltagewill be
(a)reduced by half
(b)increased by a factor of 2
(c) increased by a factor of 4
(d) unchanged
--- PAGE 28 ---
8.28Chapter8
ComprehensionVI
Billions of dollars are spent eachyear toreplace or prevent the corrosion and subsequent destruction
ofgovernment property such asships,bridges,and metallic piping.Corrosionmayresult in decreased
susceptible to corrosion. One part of the hull acts as the anode,while another acts as the cathode,the
iron hull itself connects the two parts, completing the circuit.Part of the metal hull begins oxidizing
to Fe$2$ in the presence of HO and O, the reduction reaction proceeds with an Er =1.229 V. It is
there thatFe$2$ionsmigratefrom the anode and arefurther oxidized as
2Fe$2$t(aq) + O(g) +(6 + x)H,O() $/to$Fe,O;xH,0(s) +4H,O*(aq)
2
TheFeOorrust,ormed is onlyasmall part of theproblem.Pitting,orlossofsolidmetal,weaen
the structure of the hull andmaylead to serious damage and destruction of the hull.Using a sacrificial
anode often successfully prevents corrosion.This metal,more easily oxidized than iron.Although the
sacrificialanode eventuallycorrodes aswell,the cost ofreplacingit isfarless than thatofreplacing
the iron hull.
Half-reaction
Potential
O+4H*+4e2HO
E$/circ$=1.229 V
Fe$2$+e$2$$/to$Fe$2$+
E$/circ$=0.771 V
Cu*+e$2$$/to$Cu
E$/circ$=0.521V
O+2HO+4e$/to$4OHE$/circ$=0.401V
Cu$2$+2e$2$$/to$Cu
E$/circ$=0.342 V
Pb$2$+ 2c"$/to$ Pb
E$/circ$=-0.126V
Fe$2$*+ 2e" $/to$ Fe
E$/circ$=-0.447 V
Mn$2$++ 2e$/to$ Mn
E$/circ$=-1.185 V
Mg+ 2e$/to$ Mg
E$/circ$ = -2.372 V
19.Which of the following reactions may
How many grams of Fe(s) are lost in
occur at the anode of the iron hull?
1.0 h?
(a)Fe$2$(aq)$/to$Fe$2$(aq)+e
(a) 0.261 g
(b) 0.522 g
(b)Fe(s)$/to$Fe$2$+(aq)+2e
(c) 1.044 g
(d) 0.783 g
(c)O2(g)+4H,O*(aq)$/to$6HO(1)+4e
23.A ship that has sunk to the bottom of
\underline{\hspace{1cm}}7+(3)$2$H+(be)-HO7$/leftarrow$(DO$2$H7(p)
the ocean may still exhibit corrosion. One
20.What is the standard voltage generated by
differenceinthecorrosionreactionbetween
thecorrosionreactionontheironhull of
ship above water and one that has sunk of
a ship?
thebottom of the ocean maybe that
(a) -1.676 V
(b)-0.782V
(a) the oxidation reaction of the sunken
(c) 1.782 V
(d) 1.676 V
shipdoes not include HO as a
reactant.
21.Which of the following materials could
(b) the oxidation reaction of the sunken
not serve as a sacrificial anode for lead?
shipdoes notinclude Oasareactant.
(a) Copper
(b)Iron
(c)the reduction reaction of the sunken
(c)Magnesium
(d) Manganese
shipdoes notincludeHOasa
reactant.
22.An ammeter reading shows a current of
(d)thereductionreaction of the sunken
0.50 A running through the iron hull
ship does not include Oas a reactant.
--- PAGE 29 ---
Electrochemistry8.29
ComprehensionVll
A current of 15.0A is employed to plate nickel in a NiSOsolution.Both Ni andH,are formed at
the cathode.The current efficiency with respect to formation of Ni is 60/%.The density ofnickel
=8.9 g/ml. (Ni =58.7)
24.Howmuchofnickelisplatedon the
26.Whatvolume ofH,at 0$/circ$C and1 atm is
cathode per hour?
formed per hour?
(a) 16.43 g
(b) 32.85 g
(a) 6.271
(b) 3.761
(c)19.7 g
(d) 9.85 g
(c) 2.51
(d) 5.011
25.What is the thickness of the plating if the
27.At the end of the electrolysis, how many
cathode consists of a sheet of metal4.0cm,
grams of the gaseous product appear at
which is to be coated on both sides?
the anode?
(a) 13.8mm
(b)27.6 mm
(a) 4.48 g
(b) 1.79 g
(c) 6.9 mm
(d) 23.0 mm
(c) 2.69 g
(d) 7.46 g
ComprehensionVIlI
Tollen's reagent is used for the detection of aldehyde when a solution ofAgNOis added to glucose
withNHOH then gluconic acid is formed.
A80=40+
CHO+HO$/to$CHO,(Gluconicacid)+2H++2e;E$/circ$x=-0.05V
Ag(NH)+e$/to$Ag(s)+2 NH;E$/circ$=0.337 V
[Use2.303RTIF=0.0592andFIRT=38.92 at298K]
28.2Ag*+CHO. +HO$/to$2Ag(S)
30.Ammoniaisalwaysaddedinthis
+CH,O,+ 2H*. Find InKof this reaction.
reaction.Which of thefollowingmust be
(a) 66.13
(b) 58.38
incorrect?
(c) 28.30
(d) 46.29
(a) NH, combines with Ag* to form a
complex
29.Whenammoniais added tothesolution,
(b)Ag(NH)*is a stronger oxidizing
pH is raised to 1 L. Which half-cell
agent than Ag*
reaction is affected by pH and by how
(c)In absence of ammonia, silver salt of
much?
gluconic acid is formed
(a)Eoxx will increase by a factor of 0.65
fromEd
(d)NH,has affected the standard
reduction potential of glucose-
(b)Eox will decrease by a factor of 0.65
gluconic acidelectrode
fromE$/circ$
(c)Er will increase by a factor of 0.65
from Eord
(d)Ered will decrease by a factor of 0.65
g wo
--- PAGE 30 ---
8.30Chapter8
ComprehensionIx
The standard potentialof thefollowing cell is 0.23V at 15$/circ$C and 0.21V at 35$/circ$C.
Pt(s)|H(g)|HCI(aq)|AgCI(s)|Ag(s)
Given: The standard reduction potential of the Ag*(aq)|Ag(s) couple is 0.80 V at 25$/circ$C.
H$/circ$and$/Delta$S$/circ$for the cell reaction remain unchangedin therange 15$/circ$C to35$/circ$C.
31.The overall cell reaction is
33.$/Delta$S$/circ$ for the cell reaction per mol of
(b)H$/leftarrow$(S)I7+()H ()
AgCl is
+ 2Ag(s)
(a)-193 J/K
(b)-96.5 J/K
(b) H(g)+ 2Ag(aq)$/to$ 2H"(aq) + 2Ag(s)
(c) +96.5 J/K
(d) +193 J/K
(c)Ag*(aq)+CI(aq)$/to$AgCl(s)
34.The solubility of AgClin water at 25$/circ$Cis
(d)H(g)+Cl(g)$/to$2HCI(aq)
(2.303RT/F=0.058)
32.$/Delta$H$/circ$for the cellreaction per molof
(a) 10-0 M
AgClis
(b) 10-12 M
(a)-99,974 J
(b)-49,987J
(c)10-M
(c) +99,974 J
(d)+49,987 J
(d) 10-M
ComprehensionX
For thereaction:Ag*(aq)+CI(aq)AgCl(s).Given
Species
G$2$(kJ / mol) at 25$/circ$$/circ/text(C)$
Ag*(aq)
+77
Cr(aq)
129
AgC1 (s)
109
Ag*+e$/to$Ag;E$/circ$=0.80V
90=Z+
35.Thecellrepresentationsoftheabove
37.A quantity of6.539$/times$10-$2$g ofmetallic
reaction and E are
zinc is added to 100 ml of saturated
(a)Ag(s)|AgCl(s) | Cl(aq) Il Ag*(aq)|
solutionofAgCl.The value of
Ag(s); Ea=0.59 V
[Zn$2$+]
log
is (Zn = 65.39)
(b) Ag(s) | Ag*(aq) I CI (aq) | AgCI(s) |
[Ag+]
Ag(s); Ea=0.59 V
(a) 5.288
(b) 52.88
(c) Ag(s) | CI(aq) Il Ag*(aq)| Ag; E =
(c) 528.8
(d) 26.44
0.59V
(d) Ag(s)| Ag*(aq) II CI (aq) | AgCl(s)|
38.How many moles of Ag will
Ag(s); Ea = 1.18 V
precipitated in the above reaction?
(a) 10-3
36.The solubility product of AgClat 298 K is
(b) 2$/times$10-3
(a)2$/times$10-10
(b)10-10
(c)10-6
(c) e-10
(d) 10-s
(d) 10-s
--- PAGE 31 ---
Electrochemistry8.31
ComprehensionX1
A sample of water from a large swimming poolhas a resistance of10,000 Q at 25$/circ$C,when placed
in a conductivity cell.When filled with 0.02 M-KCl solution, the cell has a resistance of 100 Q at
25$/circ$C.An amount of 585g of NaCl was dissolved in the pool,which was thoroughly stirred.A sample
of this solutiongave a resistance of 80ooQ.The molar conductivity of NaCl at this concentration is
125 Q-' cmmol' and molar conductivity of 0.02 M - KCl is 200 Q' cm$2$mol'.
39.Cell constant of theconductivity cellis
41.Volume (in L) of water in the poolis
(a) 4 cm-1
(a) 1,25,000
(b) 0.4 cm-1
(b) 12,500
(c) 40 cm-1
(c)1250
(d) 0.04 cm-1
(d)125
40.Conductivity ofwater(inQ-cm-)is
(a) 0.4
(b) 0.04
(c)0.0004
(d) 0.00004
ComprehensionXll
Redox reactions play a pivoted role in chemistry and biology.The values of standard redox potential
(E$/circ$)of two half-cell reactions decide whichway thereaction is expected to proceed.A simple example
is Daniel cell in which zinc goes into solution and copper gets deposited. Given below are a set of
half-cell reactions(acidic medium) along with their E$/circ$(Vwith respect to normal hydrogen electrode)
values.Using this data,obtain thecorrect explanations to questionsbelow:
I+2e$/to$2I
E$/circ$=0.54
Cl,+2e$2$$/to$2Cr
E$/circ$=1.36
Mn*+e"$/to$Mn$2$+
E$/circ$=1.50
E$/circ$=0.77
O+4H++4e$/to$2HO
E$/circ$=1.23
42.Among thefollowing,identify the correct
43.WhileFe't is stable,Mn$2$t is not stable in
statement?
acid solutionbecause
(a) Chloride ion is oxidized by O2
(a)Ooxidizes Mn$2$+to Mn3+
(b) Fe$2$+ is oxidized by iodine
(b) O2oxidizes both Mn2+ to Mn$2$+and
(c)Iodideion is oxidized by chlorine
Fe$2$+to Fe3+
(d)Mn$2$is oxidized by chlorine
(c)Fe$2$+oxidizes HO to O2
(d) Mn+ oxidizes HO to O2
--- PAGE 32 ---
8.32$/cdot$Chapter8
ComprehensionXil
The concentration of potassium ions inside abiological cell is at least 20 times higher than the outside.
The resulting potential difference across the cellis important in several processes such as transmission
of nerve impulses and maintaining the ion balance.A simple model for such a concentration cell
involving ametal"M'is
M(s) | M*(aq; 0.05 molar) II M*(aq; 1 molar)| M(s)
For the above electrolytic cell,the magnitude of the cell potential,|Ecel = 70 m V.
44.For the above cell
45.If the0.05molar solutionof M'is
(a) Ec<0;$/Delta$G>0
replacedby a 0.0025molarM*solution,
(b) Ec>0;$/Delta$G<0
then themagnitudeof the cellpotential
(c)Ec<0;G$/circ$>0
would be
(d) Ec>0;$/Delta$G$/circ$<0
(a)35 mV
(b)70 mV
(c)140 mV
(d)700 mV
ComprehensionXIv
Theelectrochemical cell shownbelowis aconcentration cell.
M(s)|M$2$*(saturated solution of a sparingly soluble salt,MX) II M$2$*(0.001 mol dm)| M(s)
The EMFof thecell depends onthedifferenceinconcentrations of Mions at thetwoelectrodes.The
EMF of the cell at 298 K is 0.059 V.
46.The value of $/Delta$G (kJ mol-') for the given
47.The solubility product (K;mol$2$dm-3) of
cell is (take 1 F=96500 C mol-)
MX, at 298K based on the information
(a)-5.7
availablefor the given concentration cell
(b) 5.7
is (take 2.303 $/times$ R $/times$ 298/F = 0.059)
(c) 11.4
(a) 1$/times$10-15
(b) 4$/times$10-15
(d) -11.4
(c)1$/times$10-12
(d) 4$/times$ 10-12
Comprehensionxv
Conductance measurements are frequently employed tofind the endpoints of acid-base and other
titrations.The principleinvolved is that electrical conductance of a solution depends upon thenumber
and mobility of ions.In the conductometric titrations,the conductance of the resulting solution is
measured at different stages on adding some volume of the standard solution and then a graph is
plotted from the experimentalobservations toget the endpoint.Conductometric titrations have
several advantages.Coloured solutions,titration ofweak acid and weak base,etc.cannot be titrated
by normal methods,but can be titrated successfully by this method.Further,no special care is needed
in titration because the end point is determined graphically.
--- PAGE 33 ---
Electrochemistry$/cdot$8.33
48.Which of the following graph truly
(c)
represents the titration of HCl solution
againstNaOHsolution?
(a)
Cond.
VNaOH
Cond.
(d)
VNaOH
(b)
Cond.
Cond.
VNaOH
VNaOH
50.Which of the following graph truly
(c)
represents the titration of HCl solution
againstNHOH solution?
Cond.
(a)
VNaOH
Cond.
(d)
VNHOH
Cond.
(b)
VNaOH
Cond.
49.Which of the following graph truly
VNHOH
represents the titration of CHCOOH
solution against NaOH solution?
(c)
(a)
Cond.
Cond.
VNHOH
VNaOH
(b)
(d)
Cond.
Cond.
VNaOH
VNHOH-
--- PAGE 34 ---
8.34$/cdot$Chapter8
51.Which of the following graph truly
(c)
represents the titrationof a solution
containingamixture of HCland
CHCOOHagainstNaOHsolution?
Cond.
(a)
VKCI-
Cond.
(d)
VNaOH
Cond.
(b)
VKCI
Cond.
\end{problembox}

\begin{problembox}
\textbf{Problem 58} \hfill \textbf{[Pearson Advanced Electrochemistry Q53]}
\label{q:ext-pearson-58}

Which of the following graph truly
VNaOH
represents the titration of CH,COOH
(c)
solution against NHOH solution?
(a)
Cond.
Cond.
VNaOH
(d)
VNHOH
(b)
Cond.
Cond.
VNaOH-
52.Which of the following graph truly
VNHOH
represents the titration of AgNO,
solution againstKCl solution?
(c)
(a)
Cond.
Cond.
VNHOH
VKCI
(b)
(d)
Cond.
Cond.
VKCI---$/to$
VNHOH
--- PAGE 35 ---
Electrochemistry8.35
Section D(Assertion-Reason)
The following questions consist of two
Statement II:Metals having negative
statements.Mark
reductionpotentialhave large hydration
(a)If both statements are CORRECT,
energy.
andStatement IIis the CORRECT
6.Statement I:Specific conductance
explanation ofStatement I.
decreases withdilutionwhilemolar
(b)
Ifbothstatements areCORRECT,and
conductance increases.
Statement II is NOT the CORRECT
Statement II:On dilution,number of
explanation of Statement I.
ions per unit volume decreases but total
(c)
If Statement IisCORRECT,but
number of ions increases considerably.
Statement IIisINCORRECT.
(d)If Statement I is INCORRECT,but
7.Statement I:When acidified zinc sulphate
Statement II is CORRECT.
solutioniselectrolysedbetween zinc
Statement I:Electrolysis of CuCl,(aq)
electrodes, it is zinc that is deposited at
\end{problembox}

\begin{problembox}
\textbf{Problem 59} \hfill \textbf{[Pearson Advanced Electrochemistry Q1]}
\label{q:ext-pearson-59}

the cathode and evolution of hydrogen
gives 1 mole of Cu and 1 mole of Cl, by
the passage of suitable charge.
gas does not take place.
StatementII:Theelectrode potential
Statement II:Equal equivalents of Cu
and Cl,are formed during the passage of
ofzinc becomes more negative than
hydrogen as the overvoltage for the
same charge.
hydrogen evolution on zinc is quite large.
Statement I:Zinc metal can displace Ag
metal froma solution containing the
8.Statement I:The conductivity of
complex [Ag((CN)].
solutions of different electrolytes in the
is greater than
same solvent and at a given temperature
is same.
StatementII:The conductivitydepends
\end{problembox}

\begin{problembox}
\textbf{Problem 60} \hfill \textbf{[Pearson Advanced Electrochemistry Q3]}
\label{q:ext-pearson-60}

Statement I: In electrolysis, the quantity
on the charge and size of the ions, the
of electricity neededfor depositing1 mole
concentrations ofions andease with
of silver is different from that required for
which the ions moveunder potential
1 mole copper.
gradient.
Statement II:The atomicmasses of silver
9.Statement I: Sodium ions are discharged
and copper are different.
inpreference tohydrogenions ata
Statement II:In an electrochemical cell,
mercury cathode.
\end{problembox}

\begin{problembox}
\textbf{Problem 61} \hfill \textbf{[Pearson Advanced Electrochemistry Q4]}
\label{q:ext-pearson-61}

anode and cathode are, respectively,
Statement II:Na*is a strong reducing
negative and positive electrode.
agent in comparison to H+ ion.
Statement II:At anode,oxidation takes
10.S
StatementI:The conductivityof metals
place and at cathode reduction takes
decreaseswhilethat ofelectrolytic
place.
solutionincreases with increase in
5.Statement I:A metal having negative
temperature.
reduction potential when dipped in the
Statement II:Electrons in metals are very
solution of its own ions has a tendency to
tightly held by the nucleus and are not
pass into the solution.
free to move.
--- PAGE 36 ---
8.36$/cdot$Chapter8
11.StatementI:An electrochemicalcell can
Statement II:A salt bridge contains
be set up only when the redox reaction is
concentratedsolutionofaninert
spontaneous.
electrolyte likeKCl,KNO,KSO4,etc.,
StatementII:Areaction is spontaneous
in agar-agar.
iffree energychange at constant
14.StatementI:Duringelectrolysisof
temperature and pressure is negative.
aqueous sodium acetate solution,the
\end{problembox}

\begin{problembox}
\textbf{Problem 62} \hfill \textbf{[Pearson Advanced Electrochemistry Q12]}
\label{q:ext-pearson-62}

Statement I:In the Daniel cell,if
molar ratio of gases formed at cathode
concentrations of Cu$2$+ and Zn$2$+ ions
and anode is 1:3.
are doubled, the EMF of cell does not
Statement II:Acetate ion discharges at
change.
anode and H* ion,at cathode.
Statement II:If the concentration of ions
15.Statement I:Electrode potential of any
in contact with the metal is doubled, the
electrode will change on changing any of
electrode potential will be doubled.
its intensive properties.
13.S
Statement I: KCl, NaCl, NHCl, etc.,
Statement II:Any intensive property of
cannot be used in the salt bridge of a cell
system changes on changing any of its
containing silver.
properties,whether intensive or extensive.
Section E(Column Match)
1.Match ColumnIwith Column II
Column I
Column II (Electrolysis product using inert
electrodes)
(A)Dilute solution of HCI
(P) O, evolved at anode
(B)Dilute solution of NaCl
(Q)H, evolved at cathode
(C)Concentrated solution of NaCl
(R) Cl,evolved at anode
(D)AgNO, solution
(S)Ag deposited at cathode
\end{problembox}

\begin{problembox}
\textbf{Problem 63} \hfill \textbf{[Pearson Advanced Electrochemistry Q2]}
\label{q:ext-pearson-63}

The variation in conductivity of these reactions is given in List II. Match List I with List II.
List IXY
List II
(A) (CH)N+CH,COOH
(P) Conductivity decreases and then increases
(B) KI(0.1 M) + AgNO,(0.01 M)
(Q) Conductivity decreases and then does not
change much
(C) CH,COOH+KOH
(R) Conductivity increases and then does not
change much
(D)NaOH+HI
(S)Conductivity doesnot changemuch and then
increases
--- PAGE 37 ---
Electrochemistry$/cdot$8.37
3.The standard reduction potential data at 25$/circ$C is given below:
E$/circ$ (Fe$2$, Fe$2$t)= +0.77 V
E$/circ$(Fe$2$+, Fe)= -0.44 V
E$/circ$(Cu$2$,Cu)=+0.34 V
E$/circ$(Cu*,Cu)=+0.52V
E$/circ$[O2(g)+ 4H*+4e$/to$ 2HO]=+1.23 V
E$/circ$[O2(g)+ 2H,O +4e$/to$ 4OH]= +0.40 V
E$/circ$(Cr$2$t,Cr)=-0.74 V
E$/circ$(Cr$2$+,Cr)=-0.91 V
Match E$/circ$of the redox pair in List I with the values given in List II.
List I
List II
(A) E$/circ$(Fe$2$t,Fe)
(P) -0.18 V
(B) E$/circ$(4HO$/to$4H*+4 OH')
(Q)-0.4 V
(C)E$/circ$(Cu$2$++Cu$/to$2Cu)
(R)-0.04 V
(D) E$/circ$(Cr$2$+, Cr$2$+)
(S) -0.83 V
\end{problembox}

\begin{problembox}
\textbf{Problem 64} \hfill \textbf{[Pearson Advanced Electrochemistry Q4]}
\label{q:ext-pearson-64}

Match the following
ColumnI
Column II
(A) Concentration cell
(P) Ag| AgCI|CI II Ag*| Ag
(B) Spontaneous cell reaction
(Q) Ag|AgCI| CI  Br|AgBr|Ag
(C) Non-spontaneous cell reaction
(R)Ag|Ag*(0.1M)IAg*(1.0 M)|Ag
(S)Ag|AgC1|CI (0.1 M)IICI(1.0M)|AgC1|Ag
Section F(Subjective)
Single-digit Integer Type
1.Percentageofanilinehydrochloride
potential for the above electrode becomes
hydrolysed in its M/40 solution at 25$/circ$C
0.31 V, is (2.303RT/F = 0.06)
3.The reduction potential at 25$/circ$C for
2.303RT/F= 0.06)
2.The standard reduction potential for
=-0.44 V and E
=-0.04V,theratio
\end{problembox}

\begin{problembox}
\textbf{Problem 65} \hfill \textbf{[Pearson Advanced Electrochemistry Q2]}
\label{q:ext-pearson-65}

Fe'Fe
Cu$2$|Cuelectrodeis+0.34V.The solubility
of molar concentrations of Fe$2$+ to Fe3+
product of Cu(OH), is 1.0 $/times$ 10-19. The
ions in solution is (2.303RT/F = 0.06,
pH of solution at which the reduction
log 5 = 0.7)
--- PAGE 38 ---
8.38Chapter8
4.If [Fe$2$] at equilibrium, when potassium
\end{problembox}

\begin{problembox}
\textbf{Problem 66} \hfill \textbf{[Pearson Advanced Electrochemistry Q10]}
\label{q:ext-pearson-66}

A test for complete removal of Cu2+
iodide is added toa solution ofFe+
ions from a solution of Cu$2$ is to add
initially at 0.50 M until [I]= 1.0 M,
NH,(aq).A blue colour signifies the
is x $/times$ 10- M, the value of x is (Given
formationofcomplex[Cu(NH)]
=0.53V,
having K = 1.1 x 103 and thus confirms
2.303RT/F=0.06)
I|I
the presence of Cu$2$ in solution. 250 ml
of 0.1 M -CuSO4is electrolysed by
5.If it is desired to construct the following
passing a current of 5 A for 1351 s.
galvanic cell:
After passage of this charge,sufficient
quantity of NH, is added to electrolysed
Ag(s)|Ag'(saturatedAgI)IlAg*
solution maintaining[NH)]=0.10 M.
(saturated AgCl,x$/times$10-4M CI)|Ag(s)
If [Cu(NH)]$2$+ is detectable up to its
to have Ecen = 0.102 V, what should be
concentration as low as 1 x io- M,
the value of x to get [Cl],which must
wouldablue colour be shownby the
be present in the cathodic half-cell to
electrolysedsolutiononadditionof
achievethedesiredEMF.GivenKof
NH. Mark "1', if the answer is 'yes'and
AgCl and AgI are 1.8 x 10-1 and'8.1
mark /textasciicircum()2',if the answer is/textasciicircum()no'.
x10-",respectively.(2.303RT/F =0.06)
6.An alloy of lead (valency = 2)-thallium
11.By passing a certain amount of charge
(valency=1)containing 70/%Pb and30/%
through NaCl solution,9.08Lof chlorine
T1,byweight,canbeelectroplatedontoa
gas were liberated at STP. When the same
cathode from aperchloric acid solution.
amount of charge is passed through a
How many hours would be required to
nitratesolutionofmetalM,52.8gofthe
deposit 5.0 g of this alloy at a current of
metal was deposited.If the specific heat
1.10 A?(Pb = 208,T1= 204)
of metalis 0.032 Cal/$/circ$C-g,the valency of
metal is
7.Iridium was plated from a solution
containing IrClfor 2.0 h with a current
\end{problembox}

\begin{problembox}
\textbf{Problem 67} \hfill \textbf{[Pearson Advanced Electrochemistry Q12]}
\label{q:ext-pearson-67}

The electrode reactions for charging of a
of 0.075 A.The Iridium deposited on the
lead storage battery are:
cathode weighed 0.36 g.If the oxidation
state of Ir in IrClis x, then the value of
PbSO4+ 2e $/to$Pb + SO$2$-
(x + y) is (Ir = 192)
H+-OS+$2$09d$/leftarrow$0"H7+"0S9d
8.The electrolysis of cold sodium chloride
+ 2e
solution produces sodium hypochlorite
by reacting NaOH and Cl, thoroughly
The electrolyte in the battery is an aqueous
How long (in days) will a cell operate to
solution of sulphuric acid.Beforecharging,
produce 10 L of 7.45/% (by mass) solution
the specific gravity of the liquid was found
of NaClO if the cell current is 2.5 A?
Assume that the density of solution is
to be 1.10 (16/% HSO4 by wt.).After
1.0 g/ml.
965
charging for
h, the specific gravity
9
9.A volume of 500 ml of 0.1 M-CuSO
of the liquid was found to be 1.42 (40/%
solution is electrolysed for 5 min at a
HSOby weight).If the battery contained
current of 0.161 A. If Cu is produced at
2 L of the liquid and the volume remains
one electrode and oxygen at the other, the
constant during charge,the average current
approximate pH of the final solution is
(in A) used for charging the battery is
--- PAGE 39 ---
Electrochemistry8.39
13.A dilute solution of KCl was placed
14.In the refining of silver by electrolytic
between two platinum electrodes, 12 cm
method, what will be the mass of 72.8 g
apart, across which a potential of 1.93
silver anode (60/% pure,by weight), if
volts was applied. How far (in cm) would
9.65 A current is passed for 1 h?(Ag=108)
the K+ion move in 20 h at 25$/circ$C?Ionic
\end{problembox}

\begin{problembox}
\textbf{Problem 68} \hfill \textbf{[Pearson Advanced Electrochemistry Q15]}
\label{q:ext-pearson-68}

The pH of 100 L of concentrated KCl
conductance ofK*ion at infinite dilution
at 25$/circ$C is 7.5$/times$10-$2$mho m$2$mol-'.
solutionaftertheelectrolysisfor10s
using 9.65 A at 298 K is
Four-digit Integer Type
1.The specific conductivity of a saturated
5.An alloy weighing 2.70 mg of Pb-Ag was
solution of AgCl is 2.80 $/times$ 10mho m-1
dissolved in desired amount of HNO
andvolume was made 250 ml.A silver
electrode was dipped in solution and Erll
mol-' and /textasciicircum()= 7.81 mho m$2$ mol-', the
of the cell:
solubility of silver chloride (in order of
10-g1 ) at 25$/circ$C, is
Pt, H(1 bar)| H'(1 M) I|Ag*|Ag
2.The electrode potential (in millivolts) of
was 0.50 V at 298 K.The percentage of
2Ag(s) + S$2$(aq)$/to$ Ag2S(s) + 2ein a
lead in alloy is (Given: EAg*IAg= 0.80 V,
solution buffered at pH =3 and which
Ag=108,2.303 RT/F=0.06)
is also saturated with 0.1 M -HS,is
(Given: for HS, Ka = 10-; Ka = 10-13
6.Given E$/circ$= 0.08 V for Fe$2$(cyt b)|Fe2+
Kof AgS=4x10-s, E$/circ$Agg=0.80 V,
(cyt b)couple and E$/circ$=0.20 V for
log 2 =0.3,2.303RT/F = 0.06)
Fe$2$(cyb c)|Fe$2$(cytc,)couple,where E$/circ$
represents the standard state reduction
3.An excess of liquid mercury is added to
potentials at pH = 7.0 at 25$/circ$C and cyt
an acidified solution of 10-$2$ M - Fe3.
is an abbreviation for cytochromes.The
It is found that 10/% of Fe$2$t remains at
value ofKforthereaction
equilibrium at25$/circ$C.The value ofE
Hg$2$(Hg
(in mV), assuming that the only reaction
Fe$2$ (cyt c,) +Fe$2$+ (cyt b)= Fe$2$ (cyt
that occurs is: 2Hg + 2Fe$/to$ Hg2+
c) +Fe$2$+(cyt b)
+ 2Fe*, is (Given: Ep*iFea = 0.7724 V,
is (2.303RT/F = 0.06)
log2 = 0.3,1og3 =0.48, 2.303 RT/F = 0.06)
7.Calculate potential (in mV) of the cell:
4.Determine potential (in mV) of the cell:
Cu|Mn(s)|MnCl(0.001 M),HCI(0.01
Pt|Fe$2$,Fe3CrO$2$, Cr$2$, H+|Pt in
M) O2(0.25 bar) Pt |Cu.Given,
which [Fe$2$] = 0.75 M, [Fe$2$] = 0.75 M,
E$/circ$=-1.i85 V for the Mn$2$+|Mn couple
[CrO$2$] = 2 M, [Cr$2$*] = 4 M and [H ]
and 1.229 V for the O2|HO,H couple.
=1 M. Given:
(2.303RT/F =0.06,log 2= 0.3)
Fe$2$+e$2$$/to$Fe$2$;E$/circ$=0.77 V
8.The temperature coefficient of the cell: Zn
14 H++6e +CrO$2$-$/to$ 2Cr$2$++7HO;
|ZnCl,| AgCl(s)|Ag is + 0.001 V/K. The
E$/circ$ = 1.35 V
entropy change (in J/K)accompanying
2.303RT/F=0.06,log 2=0.3,log 3=0.48
this reaction, at 298 K is
--- PAGE 40 ---
8.40$/cdot$Chapter8
9.For the reaction:4Al(s)+3O2(g)+6HO(l)
solution of Au(NO);and a current of
+4OH(aq)= 4Al(OH)(aq);Ea
2.4 A is applied. The time (in seconds)
=2.5V.If Gfor HO(l)andOH(aq)are
required for the electroplating to be
-280.0 and -156.25 kJ/mol, respectively,
completed, assuming that the layer of
the magnitude of $/Delta$G for Al(OH)4 (aq)
gold builds up evenly, is (Au = 197)
(in kJ/mol) is
\end{problembox}

\begin{problembox}
\textbf{Problem 69} \hfill \textbf{[Pearson Advanced Electrochemistry Q13]}
\label{q:ext-pearson-69}

Wehave takenasaturatedsolution
10.In a Zn-MnO,Cell, the anode is made
of AgBr. Ks, of AgBr is 12 $/times$ 10-14. If
up of Zn and cathode of carbon rod
10-moles ofAgNOis added to 1 litre of
surroundedbyamixture ofMnO2
this solution, find conductivity (specific
Carbon, NHCl and ZnCl, in aqueous
conductance)of this solution in terms of
base.If 8.7 g MnOis present in cathodic
10-7S$/cdot$m- units.Given  (Ag')=6$/times$10-3
compartment,how many days the dry
Sm$2$mol,A(Br)=8$/times$10-Sm$2$mol-,
cell will continue to give a current of
$/circ$(NO)=7$/times$10-3S$/cdot$m$2$mol-.
3.99 $/times$ 10-3 A?
\end{problembox}

\begin{problembox}
\textbf{Problem 70} \hfill \textbf{[Pearson Advanced Electrochemistry Q14]}
\label{q:ext-pearson-70}

On
passing
electricity
through
11.An amount of 19 g of molten SnCl
nitrobenzene solution,it is converted into
iselectrolysedfor sometime.Inert
azobenzene. The mass of azobenzene (in
electrodes areused.1.19g of tin is
mg) produced, if the same quantity of
deposited at the cathode.No substance
electricity produces oxygen just sufficient
is lost during electrolysis. If the ratio
to burn 96 g of fullerene (C), is
of the masses of SnCl, and SnCl4 after
15.The conductivity of saturated solution
electrolysis is 261:x,the value of x is
(Sn = 119)
of sparingly soluble salt, Ba;(PO),  is
1.2$/times$10-$/Omega$-cm-.Thelimiting
\end{problembox}

\begin{problembox}
\textbf{Problem 71} \hfill \textbf{[Pearson Advanced Electrochemistry Q12]}
\label{q:ext-pearson-71}

An object whose surface area is 80 cm$2$ is
equivalent conductances ofBacl,KPO
to be plated with an even layer of gold
and KC1are 160,140 and 100 $/Omega$cm$2$
8.0 $/times$ 104 cm thick. The density of gold
eq',respectively. The k of Bas(PO4)2
is 19.7 g/cm$2$.The object is placed in a
(in the order of 10-25)is
AnswerKeys-ExerciseIl
SectionA(onlyoneCorrect)
\end{problembox}

\begin{problembox}
\textbf{Problem 72} \hfill \textbf{[Pearson Advanced Electrochemistry Q80]}
\label{q:ext-pearson-72}

(b)
SectionB(OneorMorethanoneCorrect)
\end{problembox}

\begin{problembox}
\textbf{Problem 73} \hfill \textbf{[Pearson Advanced Electrochemistry Q15]}
\label{q:ext-pearson-73}

(a), (b), (c), (d)
--- PAGE 41 ---
Electrochemistry$/cdot$8.41
Section C
ComprehensionI
Comprehension IX
\end{problembox}

\begin{problembox}
\textbf{Problem 74} \hfill \textbf{[Pearson Advanced Electrochemistry Q34]}
\label{q:ext-pearson-74}

(c)
ComprehensionIl
ComprehensionX
\end{problembox}

\begin{problembox}
\textbf{Problem 75} \hfill \textbf{[Pearson Advanced Electrochemistry Q38]}
\label{q:ext-pearson-75}

(c)
Comprehension Ill
Comprehension XI
\end{problembox}

\begin{problembox}
\textbf{Problem 76} \hfill \textbf{[Pearson Advanced Electrochemistry Q41]}
\label{q:ext-pearson-76}

(a)
ComprehensionIV
Comprehension XlI
\end{problembox}

\begin{problembox}
\textbf{Problem 77} \hfill \textbf{[Pearson Advanced Electrochemistry Q43]}
\label{q:ext-pearson-77}

(d)
ComprehensionV
Comprehension XI
\end{problembox}

\begin{problembox}
\textbf{Problem 78} \hfill \textbf{[Pearson Advanced Electrochemistry Q44]}
\label{q:ext-pearson-78}

(b)45.(c)
Comprehension VI
Comprehension XIV
\end{problembox}

\begin{problembox}
\textbf{Problem 79} \hfill \textbf{[Pearson Advanced Electrochemistry Q47]}
\label{q:ext-pearson-79}

(b)
Comprehension VII
Comprehension XV
\end{problembox}

\begin{problembox}
\textbf{Problem 80} \hfill \textbf{[Pearson Advanced Electrochemistry Q30]}
\label{q:ext-pearson-80}

(d)
SectionD(Assertion-Reason)
\end{problembox}

\begin{problembox}
\textbf{Problem 81} \hfill \textbf{[Pearson Advanced Electrochemistry Q15]}
\label{q:ext-pearson-81}

(c)
SectionE(Column Match)
\end{problembox}

\begin{problembox}
\textbf{Problem 82} \hfill \textbf{[Pearson Advanced Electrochemistry Q1]}
\label{q:ext-pearson-82}

A$/to$P,Q;B$/to$P,Q;C$/to$Q,R;D$/to$P,S
\end{problembox}

\begin{problembox}
\textbf{Problem 83} \hfill \textbf{[Pearson Advanced Electrochemistry Q2]}
\label{q:ext-pearson-83}

A$/to$R;B$/to$S;C$/to$Q;D$/to$P
3.A$/to$R;B$/to$S;C$/to$P;D$/to$Q
\end{problembox}

\begin{problembox}
\textbf{Problem 84} \hfill \textbf{[Pearson Advanced Electrochemistry Q4]}
\label{q:ext-pearson-84}

A$/to$R,S;B$/to$P,R; C$/to$Q,S
Section F(Subjective)
Single-digitIntegerType
\end{problembox}

\begin{problembox}
\textbf{Problem 85} \hfill \textbf{[Pearson Advanced Electrochemistry Q15]}
\label{q:ext-pearson-85}

(9)
Four-digit Integer Type
\end{problembox}

\begin{problembox}
\textbf{Problem 86} \hfill \textbf{[Pearson Advanced Electrochemistry Q5]}
\label{q:ext-pearson-86}

(0090)
6.(0100)
7.(2453)
\end{problembox}